\documentclass[11pt]{article}
\usepackage{lmodern}

% ---------- Layout ----------
\usepackage[letterpaper,margin=1in]{geometry}
\usepackage{setspace}
\setstretch{1.05}

% ---------- Math & text ----------
\usepackage{amsmath,amssymb}
\usepackage{microtype}
\usepackage{longtable}
\usepackage{booktabs}
\usepackage{pdflscape}

% ---------- Figures & tables ----------
\usepackage{graphicx}
\usepackage{float}
\usepackage{array}
\usepackage{caption}
\usepackage{placeins}

% ---------- Links ----------
\usepackage[colorlinks=true,linkcolor=blue,citecolor=blue,urlcolor=blue]{hyperref}

% ---------- Bibliography ----------
\usepackage[authoryear]{natbib}

% ---------- Notation macros (00S Symbology and Notation Lock) ----------
% The notation lock is binding but had no mechanical enforcement: every symbol was
% hand-typed at every occurrence, so nothing prevented the retired \tau_t from
% reappearing in place of \kappa_t, or \delta from being reused for step dummies.
% Use these macros in new material; they are the single place to change a symbol.
\newcommand{\mut}{\mu_t}                        % capacity utilization Y/Y^p -- never "u"
\newcommand{\chit}{\chi_t}                      % recapitalization rate I/Pi -- never beta_t
\newcommand{\kapt}{\kappa_t}                    % capital composition k^ME - k^NRC -- replaces tau_t
\newcommand{\knrc}{k^{\text{NRC}}}              % nonresidential structures (log)
\newcommand{\kme}{k^{\text{ME}}}                % machinery and equipment (log)
\newcommand{\Ypot}{Y^{p}}                       % productive capacity
\newcommand{\mech}{q}                           % mechanization growth rate -- never "m"
\newcommand{\pkap}{p_t^{\kappa}}                % relative capital asset price P^ME/P^NRC
\newcommand{\impint}{\nu}                       % macro import intensity -- never mu
\newcommand{\deprec}{\delta}                    % net depreciation ONLY -- never step dummies
\newcommand{\discard}{z}                        % gross capital discard/depletion rate
\newcommand{\stepD}{\boldsymbol{\gamma}'\mathbf{D}_t}      % step dummies, single equation
\newcommand{\stepDsys}{\boldsymbol{\Gamma}\mathbf{D}_t}    % step dummies, system form
\newcommand{\MPF}{Mechanization Productive Frontier}

% ---------- Placeholder marker for underdeveloped analysis ----------
% Renders visibly in the PDF so gaps are legible to a reader rather than silent.
\newcommand{\PLACEHOLDER}[1]{%
  \par\medskip\noindent\fbox{\parbox{0.96\linewidth}{%
  \footnotesize\textbf{[PLACEHOLDER]}~#1}}\par\medskip}

% ---------- Title block ----------
\title{Distribution, Technique, and the Transformation of Accumulation: Reconstructing Capacity Utilization in the Center and Periphery\thanks{This paper benefited from AI-assisted editing for syntax refinement, structural suggestions, and stylistic consistency checks. All analytical arguments, historical interpretations, and substantive content remain the author's sole responsibility.}\\
\begin{large}
\textbf{WORK IN PROGRESS / PLEASE DO NOT CIRCULATE}     
\end{large}
}
\author{Diego Polanco\thanks{Doctoral Candidate in Economics, Department of Economics, University of Massachusetts Amherst.}\\
Department of Economics, University of Massachusetts Amherst\\
\texttt{dpolanconeco@umass.edu}}
\date{\today}

\begin{document}
\maketitle

\begin{abstract}
Capacity utilization has been treated either as a given macro ratio, an inherited benchmark, or an empirical residual. That leaves unresolved the problem of how accumulation is transformed into productive capacity under capitalist disequilibrium. This chapter reconstructs capacity utilization as a structurally mediated object, recovered through the identification of the transformation elasticity $\theta$, where distribution mediates the conversion of accumulation into productive capacity. Applying this framework comparative-econometrically to the United States (1947–2024) and Chile (1920–2024), we demonstrate that the transformation elasticity is endogenously determined through the choice of technique and the balance of class forces. In the center, the shift from a wage-led Fordist nexus to a post-Fordist relative price channel led to institutional decoupling and overaccumulation. In the periphery, this transformation is structurally truncated by external balance constraints and balance of payments dynamics. Our findings challenge the view of utilization as a simple residual, showing it to be a key institutional marker of macro-political-economic reproduction.
\end{abstract}

\vspace{1cm}

\noindent\textbf{Keywords:} Capacity Utilization; Transformation Elasticity; Choice of Technique; Cointegration; Marx-Okishio; External Balance; Center-Periphery.

\noindent\textbf{JEL Codes:} B51; C32; E01; E22; E32; P16.

\pagebreak

\FloatBarrier
\section{The Unresolved Problem of Capacity Utilization}

\paragraph{A}
This paper examines how the rate of profit, capital accumulation, and productive capacity formation interact under conditions of distributional conflict between labor and capital.

\paragraph{B}
\PLACEHOLDER{Section 1 is unwritten. What exists above is a single orienting sentence; the section that the locked outline requires does not yet exist. The remainder of this box states what it has to do, so the gap is legible rather than silent.}

\paragraph{C}
\PLACEHOLDER{\textbf{Burden of the section.} Establish that capacity utilization is unresolved in the treatments that have come down to us, in three distinct ways, and show that these are three faces of one problem rather than three separate difficulties. First, utilization treated as a \emph{given macro ratio}: taken as observed when its denominator is not. Second, utilization treated as an \emph{inherited benchmark}: anchored to an engineering or survey norm imported from outside the accounting framework that is supposed to give it meaning. Third, utilization treated as an \emph{empirical residual}: whatever is left after output has been fitted to capital, which presupposes the very transformation relation at issue. The third case is where the filtering literature sits, and the objection to it here is structural rather than methodological --- a filter that fits output to capital imposes a unit transformation elasticity by construction, and so conflates utilization with the object that would be needed to identify it. That argument is available from this chapter's own accounting framework and requires no external authority.}

\paragraph{D}
\PLACEHOLDER{\textbf{Bottom line up front.} Per the writing calibration in force, the contribution must be visible by the end of the second paragraph: what the paper examines, how it identifies the transformation elasticity, and what the comparison of the United States and Chile establishes. The reader should not have to reach Section 3 to learn what is being claimed.}

\paragraph{E}
\PLACEHOLDER{\textbf{What must not appear here.} Long literature catalogues. The full center--periphery discussion, which belongs in Section 5. Estimator debates, which belong in Section 4 and Appendix C. Detailed methodological caveats. The section motivates a problem; it does not litigate the solution.}

\paragraph{F}
\PLACEHOLDER{\textbf{Closing-line burden (locked).} The section must arrive at: \emph{the utilization problem cannot be resolved at the level of the ratio alone; it must be reconstructed through the relation that governs productive-capacity formation.} This hands off directly to Section 2.}

\paragraph{G}
\PLACEHOLDER{\textbf{Also unresolved: the abstract's sample window.} The abstract states a United States window of 1947--2024, while Sections 4 through 6 estimate over 1929--2024 with MATHPLACEHOLDER0XYZ. The two cannot both be right. Note further that 1929--2024 spans 96 years against a reported MATHPLACEHOLDER1XYZ of 94, and that the estimation prose elsewhere gives the effective span as 1931--2024 after lag-1 conditioning and listwise deletion --- which is consistent with MATHPLACEHOLDER2XYZ and suggests the table titles, not the sample, are what needs correcting.}
\FloatBarrier
\section{The Micro-Meso Hinge: Okishio, Vidal, and Capacity Formation}

\paragraph{A}
Chapter 1 established that the macroeconomic capacity ceiling is not a neutral engineering constraint, but a historically contingent outcome conditioned by the balance of class power. Yet, identifying this distributional conditioning at the macro level leaves the underlying transmission mechanism black-boxed. How, precisely, does distributive conflict physically alter the conversion of accumulated capital into productive capacity?

\paragraph{B}
Answering this requires dismantling a hidden ontological premise of heterodox growth models: the unified firm. Whether treating the normal rate of capacity utilization ($\mu_t \equiv Y_t / Y^p_t$, where $Y_t$ is realized output and $Y^p_t$ is potential capacity) as an endogenous behavioral target or an exogenous gravitational anchor, standard frameworks presume a coherent organization capable of selecting a benchmark, enforcing it across the labor process, and sustaining it over time.

\paragraph{C}
This premise collapses when confronted with the social organization of capitalist production. In a decentralized capitalist economy, individual capitals commit fixed investments without \textit{ex-ante} coordination, confronting one another in competitive markets. Within the enterprise, production decisions encounter contradictory managerial and worker imperatives. Capacity utilization cannot be reduced to a target selected by a single optimizing entity. To theorize the political economy of productive capacity, we must replace the unified firm with two interconnected frameworks: Okishio’s theory of decentralized market anarchy, and Vidal’s theory of the internally divided firm.

\subsection{Smooth Reproduction and the Unified-Firm Premise}

\paragraph{A}
Smooth capitalist reproduction requires strict material proportionality between departments of production, a condition formally proven by Basu (2022) and defined by Okishio (2022) as the equilibrium accumulation trajectory. Bergmann et al. (2009) refer to this balanced trajectory—where realized output meets productive capacity ($\mu_t = 1$)—as the "golden path." However, this proportional benchmark serves strictly as an analytical condition of feasibility, not a historical tendency or a gravitational attractor. While a utilization rate of unity ($\mu_t = 1$) is a necessary condition for balanced reproduction, it is not sufficient; full macroeconomic reproduction also requires that the rate of accumulation matches the growth rate of the labor force and technical change. Nothing in the logic of decentralized accumulation establishes that competition will reproduce these proportions or pull the macroeconomy back toward unity.

\paragraph{B}
Because private investment decisions are made without \textit{ex-ante} coordination, real capitalist economies routinely violate the strict conditions of the golden path. Okishio (2022) demonstrates that this chronic disproportionality is not a temporary friction but the structural norm of decentralized accumulation. Treating this transient benchmark as a behavioral target relies on the flawed premise of the unified firm.

\paragraph{C}
In reality, the denominator of the utilization ratio (productive capacity $Y^p_t$) is not an engineering constant, but a contested settlement of divided management and shop-floor struggle; the numerator (realized output $Y_t$) is similarly hostage to the anarchy of demand. Smooth reproduction supplies a reference case against which this historical non-coincidence can be studied; it does not describe the path capitalist accumulation actually follows.

\subsection{Decentralized Accumulation and Market Anarchy}

\paragraph{A}
If the “golden path” of smooth reproduction is merely a benchmark of feasibility, the actual trajectory of capitalist accumulation is governed by the decentralized organization of investment. Capitalist production is coordinated through exchange, meaning that the social validation of private decisions occurs only \textit{ex-post}. Because money mediates exchange, sale and purchase are separated moments of the capital circuit (Basu, 2022). This constitutes the logic of capital as impersonal domination: capitalists make their production and investment decisions with the law of value operating \textit{“behind their backs”} \footnote{Karl Marx, \textit{Capital}, Volume I, Chapter 1, Section 4 ("The Fetishism of Commodities and the Secret thereof"), where he describes the law of value operating as an overriding law of nature when private producers exchange products.}. They are subject to market forces and value relations they neither control nor consciously direct. Okishio (2022) locates the source of macroeconomic instability precisely in this structural anarchy. Individual capitals commit to fixed investment on the basis of private profit expectations, yet the market validates those commitments only after production has begun. No firm can know in advance whether its expansion fits the investment and consumption decisions taken elsewhere. The coordination problem is not one of imperfect information surrounding an otherwise coherent plan; no economy-wide plan exists.

\paragraph{B}
Smooth reproduction requires proportional expansion between Department I, which produces means of production, and Department II, which produces consumption goods. Because private investment does not enforce this proportionality, firms extrapolate past prices, profit rates, and sales. A signal that induces expansion in one branch routinely produces excess capacity before complementary demand materializes elsewhere. Contraction works through the same cumulative mechanism: falling sales cut investment, reducing aggregate demand and validating further cuts. Prices register the resulting imbalance, but they do not prevent the sunk commitments that produced it.

\paragraph{C}
Fixed capital gives this temporal lag its structural force. Structures and machinery embody long-horizon commitments that cannot be reversed when relative prices or demand conditions shift. Once installed, an unvalidated project leaves the firm to operate below its anticipated scale, defer maintenance, hold equipment idle, or devalue the asset. These responses distribute the costs of market anarchy among firms, workers, and creditors, but they do not restore the missing \textit{ex-ante} coordination.

\paragraph{D}
Idle capacity is therefore a recurrent, structural form of decentralized accumulation rather than a temporary friction. Private investment physically builds capacity before the market can validate its social use, repeatedly generating a gap between installed productive potential ($Y^p_t$) and realized output ($Y_t$). In this framework, the capacity utilization rate ($\mu_t$) operates as the macroeconomic "exhaust valve" for structural disproportionality. Rather than adjusting investment to smoothly restore the golden path of balanced growth, the system absorbs the shock of uncoordinated accumulation through fluctuating utilization rates. The condition $\mu_t = 1$ is therefore at most a transient point of accidental proportionality, while persistent divergence records the ongoing reality of market anarchy.

\paragraph{E}
Yet Okishio’s macro framework leaves a crucial gap: it does not explain how firms translate this external anarchy into internal operating decisions. A firm does not receive a price signal and respond as a frictionless, unified unit. It processes the signal through an organization divided over control, cooperation, effort, and the allocation of resources. Vidal’s internally divided firm provides that missing mediation.

\subsection{The Internally Divided Firm}

\paragraph{A}
Vidal reconstructs Marx’s Volume I theory of the labor process for the contemporary organizational era. Marx established that capitalist management faces a constitutive contradiction: it must simultaneously organize the cooperative labor required for production and enforce the coercive control required for valorization (surplus extraction). Vidal demonstrates how this contradiction fractures the modern capitalist firm, treating it as a divided social organization rather than a single-minded production function (Vidal, 2013, 2022). Management must extract surplus value through control over labor, but production also requires workers’ cooperation, tacit knowledge, and problem-solving capacity. Standardization and surveillance secure a measurable pace of work while weakening the discretion needed to sustain complex production. The firm cannot maximize both imperatives at once because the labor process that produces value is also the site where workers contest its extraction.

\paragraph{B}
Workers respond to this contradiction by actively calibrating effort, cooperation, machine pacing, and the deployment of tacit knowledge. Capitalist management extracts potential capacity through concrete shop-floor mechanisms: increasing machine speed (speedups), deferring maintenance to hit short-term throughput targets, and enforcing rigid surveillance. Workers contest this extraction through labor turnover, collective exit, or the deliberate withdrawal of problem-solving cooperation and effort. Throughput, quality, and maintenance therefore depend on social relations that no engineering specification fixes in advance. A production target that raises machine speed by deferring maintenance may increase current output while degrading sustainable capacity. A directive that suppresses worker discretion may meet a short-term labor-control objective while worsening workflow bottlenecks. Labor effort is consequently not a fixed technical parameter applied to a pre-existing capital stock; it is an active, contested determinant of the capacity denominator itself.

\paragraph{C}
This dynamic enforces a strict wedge between \textit{engineering capacity} and \textit{practical capacity}. Engineering capacity records the theoretical maximum output of equipment under ideal technical conditions. Practical capacity records the sustainable output rate dictated by actual maintenance schedules, product-mix complexity, and the negotiated intensity of labor. When managers report capacity to national statistical agencies, they rely on conventions shaped by their organizational position and the immediate pressures of the shop floor. Aggregate capacity statistics therefore do not reveal a neutral physical ceiling; they aggregate situated, institutional judgments. For the purposes of macroeconomic modeling, the productive capacity ceiling ($Y^p_t$) is not an exogenous engineering constraint, but an endogenous settlement of workplace conflict.

\paragraph{D}
Because the valorization-labor process contradiction cannot be solved, firms settle into organizational configurations of "mediocre sufficiency" (Vidal, 2022). These are hybrid arrangements that reproduce production without resolving the underlying conflict between control and cooperation. Partial worker autonomy coexists with rigid discipline; continuous improvement programs coexist with brutal throughput pressure; idle machinery coexists with binding bottlenecks. These arrangements are not temporary frictions or departures from an otherwise unified, optimizing firm. They are the permanent, structural forms through which capitalist management physically converts accumulated capital into usable capacity.

\subsection{Distribution, Technique, and the Transformation of Accumulation}

\paragraph{A}
The choice of technique connects the divided firm to distribution, governing how accumulated capital is physically transformed into productive capacity. Capitalists adopt a new technique when it reduces costs at prevailing prices and real wages. Okishio’s theorem shows that a viable capital-using, labor-saving (CU-LS) technique raises the profit rate when the real wage remains fixed (Okishio, 1961). Marx’s account allows the same form of technical change to depress profitability when rising capital intensity and the distribution of productivity gains increase the capital advanced relative to surplus value. The Marx–Okishio threshold identifies the rate of real-wage growth that separates these outcomes (Basu, 2022). It is a distributive boundary set by the balance of class forces, not a technical constant.

\paragraph{B}
This boundary gives wage pressure a precise role in mechanization. When organized labor compels real wages to follow productivity, labor-saving techniques become more attractive while workers retain part of the resulting productivity gain. Under the institutional settlement of Atlantic Fordism—what Basu (2026) formalizes as the "strong labor" closure—this mechanism joined rapid mechanization to continuing conflict over profitability (Vidal, 2015, 2019). Wage-led mechanization was not a harmonious productivity bargain; the same distributive pressure that induced technical change contested capital’s claim on its proceeds.

\paragraph{C}
The weakening of organized labor after the Fordist crisis altered that mechanism. Under the "weak labor" Post-Fordism, wage suppression reduced the distributive pressure to replace labor, even as lower relative machinery prices and technical change in capital-goods production continued to favor mechanization. Distribution and machinery prices therefore constitute a \textbf{dual-channel mechanism} in the choice of technique. Their relative force changes across accumulation regimes, which prevents the wage share from being treated as an external control added to an otherwise autonomous technical relation.

\paragraph{D}
These channels still operate through the contradictions identified above. Firms choose techniques before demand validates the associated capacity, and they operate those techniques through labor processes that divide control from productive cooperation. Investment consequently does not translate one-for-one into a stable productive ceiling. The transformation depends on the composition of capital, the distributive conditions governing mechanization, and the organization through which installed equipment becomes usable capacity.

\paragraph{E}
Okishio and Vidal thus rule out a direct passage from accumulation to productive capacity. Decentralized investment generates disproportionality across firms and sectors; the internally divided firm transmits that disproportionality into contested operating conditions. The relevant object is therefore the aggregate \textbf{transformation elasticity ($\theta \equiv d y^p / d k$)}, defined in log-levels to ensure a dimensionless parameter, which governs how the accumulation of capital translates into potential capacity expansion under distribution and market anarchy.

\subsection{Historical Regimes of Capacity Formation in the Center and the Periphery}

\paragraph{A}
The theoretical framework developed above—combining Okishio’s market anarchy, Vidal’s divided firm, and the Marx-Okishio threshold—provides the political economy foundation for understanding how the transformation elasticity ($\theta$) shifts across distinct historical eras. Mapping the dual-channel mechanism of technical change onto the institutional history of center and peripheral economies yields three core hypotheses regarding the structural evolution of capacity formation:

\paragraph{B}
\textbf{Hypothesis 1: The Fordist Distributive-Mechanization Nexus in the Center (1947–1973)} Under the institutional settlement of Atlantic Fordism, rising wage shares acted as the primary catalyst for intensive mechanization. Capitalists responded to sustained distributive pressure by aggressively substituting machinery for labor. This wage-induced choice of technique was structurally accommodated by institutionally coupled capacity planning: extensive plant structures and intensive machinery expanded in tandem with the new technical composition of capital. This regime was structurally stabilized by the mid-1950s—leading to a key econometric break in \textbf{1956} characterized by post-Korean War defense budget stabilizers, long-horizon corporate planning, and labor union bureaucratization. These institutions locked in capacity designs but began to restrict capacity utilization flexibility, sustaining productive capacity formation until the late-Fordist profit squeeze (circa 1968–1973) exhausted the regime's mediating capacity, triggering a systemic crisis of accumulation.

\paragraph{C}
\textbf{Hypothesis 2: Post-Fordist Institutional Decoupling and Unproductive Overaccumulation in the Center (1974–Present)} Under Post-Fordism, the suppression of labor's bargaining power neutralized the distributive driver of mechanization. This transition was triggered by the structural breaks of \textbf{1974} (the Bretton Woods collapse, OPEC oil shock, and severe recession) and \textbf{1980} (the Volcker interest rate shock and Reagan's anti-labor counter-offensive). These shocks permanently depressed the baseline level of capacity validation by establishing a neoliberal regime of wage repression and prioritizing short-horizon, financialized asset management over physical capacity utilization. While aggregate capital accumulation continued, it fundamentally decoupled from structural plant expansion. In the absence of wage-induced productive innovation, the \textit{relative price channel} (the cheapening of short-horizon equipment and ICT) dominated the choice of technique. Capital was increasingly allocated to non-productive, financially reversible processes rather than long-lived productive capacities. This institutional decoupling generated a chronic capacity realization failure, resulting in a severe structural overaccumulation regime where marginal capital growth yields diminishing or negative potential capacity expansion.

\paragraph{D}
\textbf{Hypothesis 3: Peripheral Subordination and the Balance of Payments Truncation (1920–Present)} In the capitalist periphery, the transformation of accumulation into productive capacity is structurally truncated by external financial subordination and the Balance of Payments (BoP) constraint. Unlike the center, where domestic class compromise or financialization dictates the choice of technique, peripheral mechanization relies heavily on imported capital goods and technologically non-substitutable inputs. When external financing compressions bind—whether under the inward-looking state-led regimes of the mid-century or the financialized neoliberal export models of the post-Fordist era—foreign exchange scarcity generates a hard realization bottleneck. The effective profit share available for reinvestment is squeezed, blocking the materialization of accounting surplus into productive capacity. This external subordination forces a chronic peripheral accumulation wedge: nominal profitability coexists with truncated mechanization and stalled structural envelopes. Consequently, the peripheral productive ceiling is locked between domestic distributive struggles and the disciplinary regime of global financial markets, permanently destabilizing the long-run coupling between plant and machinery.

\paragraph{E}
These hypotheses establish that the productive ceiling is not a neutral engineering constraint, but a contingent outcome of decentralized accumulation, the balance of class power, and a country's geopolitical position within the world system. The immediate task is to translate these historical-institutional claims into a formal mathematical architecture. Section 3 will therefore formalize the dual-channel choice of technique, the heterogeneous composition of capital, and the peripheral foreign-exchange bottleneck to derive the regime-specific transformation elasticity ($\theta$) that governs the passage from accumulation to productive capacity.
\FloatBarrier
\section{The Analytical Framework: Formalizing the
	Transformation Elasticity}

\paragraph{3.0}
This section translates the distributive--technological
mechanism of Section~2 into a formal model that yields
estimable restrictions. \citet{CajasGuajirro2024} first
advanced the cost-minimization problem over mechanization
growth; we extend it through three subsections of
increasing complexity. Section~3.1 formalizes aggregate
capital-deepening under Leontief technology and derives
the transformation elasticity $\theta$ endogenously from
the cost-minimizing choice of technique. Section~3.2
unbundles aggregate capital into extensive nonresidential
structures and intensive machinery, and derives the
state-dependent mechanization gap $\kappa_t$ under a
dual-channel viability condition. Section~3.3 imposes the
peripheral foreign-exchange bottleneck as a binding
constraint on $\kappa_t$, generating a corner solution
that truncates $\theta$. Each subsection nests the
preceding one: setting the heterogeneity parameter to
zero recovers the baseline model; relaxing the
foreign-exchange constraint recovers the center model.
The empirical analysis bounds this multi-sector mechanism
to a single-sector aggregate, which masks internal
sectoral divergences but preserves the core transformation
logic. Section~3.4 synthesizes the theoretical claims of
the three models. Section~4 consolidates the
identification strategy and derives the cointegrating
specifications.

\subsection{The Baseline Model of Cost-Minimizing
	Mechanization (Homogeneous Capital Stock)}

\paragraph{3.1}
Actual aggregate output is represented by a Leontief
production function with fixed technical coefficients:
\begin{equation}
	Y = \min\left( A L,\; \mu B K \right)
	\label{eq:31:leontief}
\end{equation}
where $A = Y/L$ is actual labor productivity, $\mu = Y/Y^p$
is capacity utilization, and $B = Y^p/K$ is maximum capital
productivity. This production structure governs the rate of
profit through the fundamental heterodox accounting identity:
\begin{equation}
	r = \mu B (1 - \omega)
	\label{eq:31:profit_identity}
\end{equation}
where $\omega$ is the wage share. Standard macroeconomic
models treat $B$ as an exogenous technological given and $\mu$
as a behavioral constant smoothly targeted by a unified firm.
We reject both assumptions. The components of
\eqref{eq:31:profit_identity} are endogenous outcomes of
class struggle and decentralized market anarchy, as
established in Section~2. This subsection formalizes how
distributive pressure forces cost-minimizing firms to alter
their choice of technique, endogenously determining $B$.

\subsubsection{Tactical Choice of Technique and Induced
	Innovation}

\paragraph{3.2}
Firms select a mechanization growth rate $q$ to maximize the
net productivity surplus. Let $a$ denote the growth rate of
labor productivity. The effective distributive cost of
deploying mechanization growth is $q(1-\omega)$. The firm
solves:
\begin{equation}
	\max_{q}\; c(q) = a - q(1-\omega)
	\label{eq:31:objective}
\end{equation}
This objective captures the induced-innovation logic: firms
accelerate mechanization to expand labor productivity growth,
but the distributive cost of capital intensity holds them
back. A technological frontier bounds this choice:
\begin{equation}
	a = \Phi(q) = \lambda_0 + \lambda_1 q + \lambda_2 q^2
	\label{eq:31:frontier}
\end{equation}

\paragraph{3.3}
\textbf{Assumption 1.} \textit{Concavity of the
	mechanization frontier.} The frontier parameters satisfy
$\lambda_0 > 0$, $\lambda_1 > 0$, and $\lambda_2 < 0$. The
base productivity growth $\lambda_0$ sets the floor without
advanced machinery. The linear coefficient $\lambda_1$
captures the direct productivity payoff of mechanization. The
strict concavity $\lambda_2 < 0$ ensures diminishing marginal
returns to mechanization.

\paragraph{3.4}
Substituting \eqref{eq:31:frontier} into
\eqref{eq:31:objective} and solving the first-order condition
yields the optimal mechanization growth:
\begin{align}
	c(q) &= \lambda_0 + \left[\lambda_1 - (1-\omega)\right] q
	+ \lambda_2 q^2 \label{eq:31:unconstrained} \\
	\frac{dc}{dq} &= \lambda_1 - (1-\omega) + 2\lambda_2 q^*
	= 0 \label{eq:31:foc} \\
	q^*(\omega) &= \frac{(1-\omega) - \lambda_1}{2\lambda_2}
	\label{eq:31:qstar}
\end{align}
The strict concavity ($c''(q) = 2\lambda_2 < 0$) guarantees
a unique maximum.

\paragraph{3.5}
\textbf{Proposition 1.} \textit{Induced innovation.} Under
Assumption~1, the optimal mechanization growth is strictly
increasing in the wage share:
\begin{equation}
	\frac{\partial q^*}{\partial \omega}
	= -\frac{1}{2\lambda_2} > 0
	\label{eq:31:comparative_static}
\end{equation}
A rise in the wage share forces cost-minimizing firms to
substitute machinery for labor. Rising wages drive
mechanization directly through the first-order condition,
pushing optimal capital intensity growth upward.
\textit{Proof.} Immediate from differentiating
\eqref{eq:31:qstar}. $\hfill\blacksquare$

\paragraph{3.X}
\textbf{Remark 1.} The objective
\eqref{eq:31:objective} does not assert that the firm is a
coherent optimizer. Following \citet{Vidal2022}, the firm is
internally divided between the \textit{control imperative}
(valorization: cost reduction through mechanization) and the
\textit{cooperation imperative} (production: the need for
worker effort, tacit knowledge, and problem-solving
capacity). The formal model captures the cost-minimizing
pole---the Okishio viability logic---while the constraints on
the optimization problem capture the limits imposed by the
cooperation imperative and by external conditions. The firm
maximizes net productivity surplus---labor productivity
growth minus the distributive cost of mechanization---subject
to diminishing returns. The firm's actual mechanization
choice is a contested settlement between these imperatives,
not a coherent optimum. What the model demonstrates is the
cost-reducing direction in which distributive pressure pushes
the settlement, not the settlement itself.

\subsubsection{Endogenous Transformation Elasticity}

\paragraph{3.6}
Substituting the optimal technique back into the
technological frontier recovers the normal growth rate of
labor productivity under smooth reproduction:
\begin{equation}
	a^N(\omega) = \lambda_0 + \lambda_1 q^*(\omega)
	+ \lambda_2 \left[ q^*(\omega) \right]^2
	\label{eq:31:aN}
\end{equation}
This optimal technique dictates the structural trajectory of
the Leontief parameter $B$. By definition, maximum capital
productivity is $B = Y^p / K$. In dynamic growth rates, the
growth of capital productivity is:
\begin{equation}
	b \equiv g_{Y^p} - g_K = a - q
	\label{eq:31:b_def}
\end{equation}
Evaluated at the smooth reproduction benchmark:
$b^N(\omega) = a^N(\omega) - q^*(\omega)$.

\paragraph{3.7}
\textbf{Definition 1.} \textit{Transformation elasticity.}
The transformation elasticity $\theta$ is the structural
parameter that translates capital accumulation into productive
capacity growth:
\begin{equation}
	g_{Y^p} = \theta\, g_K
	\label{eq:31:theta_def}
\end{equation}
Substituting \eqref{eq:31:theta_def} into
\eqref{eq:31:b_def} reveals the link between the level
coefficient $B$ and the elasticity $\theta$:
\begin{equation}
	\theta = 1 + \frac{b}{g_K}
	\label{eq:31:theta_link}
\end{equation}

\paragraph{3.8}
To evaluate this elasticity, we anchor the optimal technique
to a baseline rate of capital accumulation. Following the
normal reproduction path, the normal accumulation rate is
driven by the normal profit rate and the normal
recapitalization share, net of depreciation:
\begin{equation}
	g_K^N = \chi^N r^N - \delta
	\label{eq:31:gKN}
\end{equation}
Expressed dynamically relative to this normal path, the
aggregate transformation elasticity becomes:
\begin{equation}
	\theta(\omega) = 1 + \frac{b^N(\omega)}{g_K^N}
	= 1 + \frac{a^N(\omega) - q^*(\omega)}{g_K^N}
	\label{eq:31:theta_omega}
\end{equation}

\paragraph{3.9}
\textbf{Proposition 2.} \textit{Endogeneity of the
	transformation elasticity.} Under Assumption~1, the
transformation elasticity $\theta$ is not an exogenous
technical constant. It emerges endogenously from the
cost-minimizing choice of technique. When distributive
pressure forces firms to alter mechanization via
\eqref{eq:31:comparative_static}, it simultaneously shifts
the structural efficiency of accumulation via
\eqref{eq:31:theta_omega}. The elasticity $\theta$ is
therefore a function of the wage share $\omega$, the frontier
parameters $\lambda_0, \lambda_1, \lambda_2$, and the normal
accumulation rate $g_K^N$.


\subsection{Heterogeneous Capital, the Mechanization Gap,
	and State-Dependent Viability}

\paragraph{3.10}
Mainstream macroeconomics treats the capital stock as a
homogeneous aggregate, abstracting from the internal
architecture of the firm. Following the theory of the
divided firm \citep{Vidal2015, Vidal2019}, we unbundle
aggregate productive capital into extensive nonresidential
structures ($K^{NRC}$) and intensive machinery ($K^{ME}$).
These components do not enter the labor process
symmetrically. Structures form the passive infrastructure
envelope---the factory floor and logistics shell. Machinery
is the active component that restructures the production of
relative surplus value \citep[Vol.~I, Ch.~15]{MarxKVI}.

\paragraph{3.11}
\textbf{Assumption 2.} \textit{Heterogeneous capital
	decomposition.} Aggregate productive capital decomposes as
$K^{cap}_t = K^{NRC}_t + K^{ME}_t$. Nonresidential
structures evolve independently of short-run distributive
conflict ($\omega$) and follow a normal accumulation
corridor governed by long-run reproduction conditions:
\begin{equation}
	g_{K^{NRC}, t}^N = \chi^N r^N_t - \delta
	\label{eq:32:gKNRC}
\end{equation}
This extensive envelope $\Gamma(K^{NRC})$ bounds the space
of tactical technique choice.

\subsubsection{The Extensive Envelope and the Intensive
	Margin}

\paragraph{3.12}
Given the fixed infrastructural envelope, the
\textbf{intensive margin} of accumulation is the tactical
choice of the mechanization gap, defined in natural logs as:
\begin{equation}
	\kappa_t = k^{ME}_t - k^{NRC}_t
	\label{eq:32:kappa_def}
\end{equation}
This gap measures mechanization intensity relative to the
extensive plant scale. The capacity identity decomposes
potential output into the extensive envelope and the
intensive capacity yield $\Psi_t$:
\begin{equation}
	y^p_t = \theta_{NRC} k^{NRC}_t
	+ \Psi_t(\kappa_t, \omega_{t-1}, p^\kappa_t)
	\label{eq:32:capacity_decomp}
\end{equation}
The intensive yield $\Psi_t$ is not a fixed engineering
constant. It is the outcome of a decentralized
cost-minimization problem governed by the regime of
accumulation, as formalized below.

\subsubsection{Cost-Surplus Maximization and the Okishio
	Viability Yield}

\paragraph{3.13}
The firm chooses the mechanization gap $\kappa_t$ to
maximize the net viability surplus, subject to the
regime-conditioned mechanization frontier:
\begin{equation}
	\Phi(\kappa_t) = \lambda_0 + \lambda_1 \kappa_t
	+ \lambda_2 \kappa_t^2
	\label{eq:32:frontier_kappa}
\end{equation}
The \textit{realized} capacity yield is governed by the
\textbf{Marx--Okishio viability condition}: a
cost-minimizing firm deploys a new technique only if it
reduces unit costs at current distributive and price
parameters \citep{Okishio1961, Basu2022}. We formalize
this as a cost-surplus maximization problem:
\begin{equation}
	\max_{\kappa_t}\; c(\kappa_t)
	= \Phi(\kappa_t) - \kappa_t \cdot MC(\omega_{t-1}, p^\kappa_t)
	\label{eq:32:objective_kappa}
\end{equation}

\paragraph{3.14}
The marginal viability cost $MC(\omega_{t-1}, p^\kappa_t)$
is determined by two channels:
\begin{enumerate}
	\item \textbf{The distributive dividend (labor saving).}
	Increasing $\kappa_t$ substitutes machinery for labor. The
	value of this displacement is proportional to the lagged
	wage share ($\omega_{t-1}$), reflecting the wage
	settlement period prior to current capacity realization.
	A higher wage share \textit{reduces} the net cost of
	mechanization.
	\item \textbf{The relative price penalty (capital cost).}
	Increasing $\kappa_t$ requires accumulating machinery
	faster than structures. The cost of this composition shift
	is proportional to the log relative price of machinery
	($p^\kappa_t = \ln P^{ME}_t - \ln P^{NRC}_t$). A higher
	relative price \textit{increases} the marginal cost.
\end{enumerate}
We specify the marginal viability cost as:
\begin{equation}
	MC(\omega_{t-1}, p^\kappa_t)
	= \gamma_0 - \gamma_1 \omega_{t-1} + \gamma_2 p^\kappa_t
	\label{eq:32:MC}
\end{equation}
where $\gamma_1 > 0$ captures the wage-savings dividend and
$\gamma_2 > 0$ captures the relative price penalty.

\paragraph{3.15}
Substituting \eqref{eq:32:frontier_kappa} and
\eqref{eq:32:MC} into \eqref{eq:32:objective_kappa} and
solving the first-order condition yields the optimal
mechanization gap:
\begin{align}
	c(\kappa_t) &= \lambda_0
	+ \left[\lambda_1 - (\gamma_0 - \gamma_1 \omega_{t-1}
	+ \gamma_2 p^\kappa_t)\right] \kappa_t
	+ \lambda_2 \kappa_t^2
	\label{eq:32:unconstrained_kappa} \\
	\frac{dc}{d\kappa_t} &= \lambda_1 - \gamma_0
	+ \gamma_1 \omega_{t-1} - \gamma_2 p^\kappa_t
	+ 2\lambda_2 \kappa_t^* = 0
	\label{eq:32:foc_kappa} \\
	\kappa_t^*(\omega_{t-1}, p^\kappa_t)
	&= \frac{(\gamma_0 - \lambda_1) - \gamma_1 \omega_{t-1}
		+ \gamma_2 p^\kappa_t}{2\lambda_2}
	\label{eq:32:kappa_star}
\end{align}

\paragraph{3.16}
\textbf{Proposition 3.} \textit{Dual-channel induced
	innovation.} Under Assumption~1 (concavity) and
Assumption~2 (heterogeneous capital), the optimal
mechanization gap satisfies:
\begin{align}
	\frac{\partial \kappa_t^*}{\partial \omega_{t-1}}
	&= -\frac{\gamma_1}{2\lambda_2} > 0
	\label{eq:32:cs_omega} \\
	\frac{\partial \kappa_t^*}{\partial p^\kappa_t}
	&= \frac{\gamma_2}{2\lambda_2} < 0
	\label{eq:32:cs_price}
\end{align}
A rising wage share forces cost-minimizing firms to widen
the mechanization gap, substituting machinery for labor.
A rising relative price of machinery suppresses the optimal
gap, as the structural cost of the composition shift
penalizes mechanization. \textit{Proof.} Immediate from
differentiating \eqref{eq:32:kappa_star}. $\hfill\blacksquare$

\subsubsection{The Marginal Capacity Payoff and the
	Econometric Bridge}

\paragraph{3.17}
The macroeconomic translation from the mechanization gap to
potential capacity is given by the \textbf{marginal
	capacity payoff} ($MP_\kappa$), the derivative of the
intensive capacity yield $\Psi_t$ with respect to $\kappa_t$.
By the envelope theorem:
\begin{align}
	MP_\kappa &= \frac{\partial \Psi_t}{\partial \kappa_t}
	= \frac{\partial \Phi}{\partial \kappa_t}
	- MC(\omega_{t-1}, p^\kappa_t)
	\label{eq:32:MP_env} \\
	&= (\lambda_1 + 2\lambda_2 \kappa_t)
	- (\gamma_0 - \gamma_1 \omega_{t-1} + \gamma_2 p^\kappa_t)
	\label{eq:32:MP_expanded} \\
	&= \underbrace{(\lambda_1 - \gamma_0
		+ 2\lambda_2 \kappa_t)}_{\text{autonomous yield}}
	+ \gamma_1 \omega_{t-1} - \gamma_2 p^\kappa_t
	\label{eq:32:MP_rearranged}
\end{align}

\paragraph{3.18}
Integrating the marginal payoff with respect to $\kappa_t$
(treating $\omega_{t-1}$ and $p^\kappa_t$ as predetermined
shifters) and defining $\theta_\kappa$,
$\theta_{\omega\kappa} \equiv \gamma_1$, and
$\theta_{\kappa p\kappa} \equiv -\gamma_2$ yields the
intensive capacity yield function:
\begin{equation}
	\Psi_t = \theta_\kappa \kappa_t
	+ \theta_{\omega\kappa} (\omega_{t-1} \kappa_t)
	+ \theta_{\kappa p\kappa} (p^\kappa_t \kappa_t)
	\label{eq:32:Psi}
\end{equation}
Substituting \eqref{eq:32:Psi} back into
\eqref{eq:32:capacity_decomp} recovers the structural
capacity equation:
\begin{equation}
	y^p_t = \theta_{NRC} k^{NRC}_t + \theta_\kappa \kappa_t
	+ \theta_{\omega\kappa} (\omega_{t-1} \kappa_t)
	+ \theta_{\kappa p\kappa} (p^\kappa_t \kappa_t)
	\label{eq:32:capacity_eq}
\end{equation}
The interaction terms are not ad-hoc econometric additions.
They are the macroeconomic footprint of the divided firm's
Okishio viability constraint. Distributive pressure and
relative price enter because they alter the marginal
capacity payoff of the mechanization gap.

\subsubsection{The Aggregate Transformation Elasticity}

\paragraph{3.19}
The aggregate transformation elasticity
($\theta^{SpecB}_t$) is the ratio of the total differentials
of potential capacity and aggregate log capital along the
historical accumulation path. The total differential of
\eqref{eq:32:capacity_eq}, holding the slow-moving state
variables fixed, separates the extensive envelope from the
intensive margin:
\begin{align}
	dy^p_t &= \theta_{NRC} dk^{NRC}_t + MP_\kappa d\kappa_t
	\label{eq:32:dy_total} \\
	&= \theta_{NRC} dk^{NRC}_t
	+ \left(\theta_\kappa + \theta_{\omega\kappa} \omega_{t-1}
	+ \theta_{\kappa p\kappa} p^\kappa_t\right) d\kappa_t
	\label{eq:32:dy_expanded}
\end{align}

\paragraph{3.20}
The variation in aggregate capital follows from the Divisia
index of the heterogeneous components. Since
$s_{ME,t} + s_{NRC,t} = 1$ and
$d\kappa_t = dk^{ME}_t - dk^{NRC}_t$:
\begin{align}
	dk^{cap}_t &\approx s_{ME,t} dk^{ME}_t
	+ s_{NRC,t} dk^{NRC}_t
	\label{eq:32:dk_divisia} \\
	&= dk^{NRC}_t + s_{ME,t} d\kappa_t
	\label{eq:32:dk_simplified}
\end{align}
Dividing \eqref{eq:32:dy_expanded} by
\eqref{eq:32:dk_simplified} yields the elasticity along the
empirical growth path:
\begin{equation}
	\theta^{SpecB}_t
	= \frac{\theta_{NRC}
		+ \left(\theta_\kappa + \theta_{\omega\kappa} \omega_{t-1}
		+ \theta_{\kappa p\kappa} p^\kappa_t\right)
		\frac{d\kappa_t}{dk^{NRC}_t}}
	{1 + s_{ME,t} \frac{d\kappa_t}{dk^{NRC}_t}}
	\label{eq:32:theta_SpecB}
\end{equation}

\paragraph{3.21}
The aggregate transformation elasticity
$\theta^{SpecB}_t$ in \eqref{eq:32:theta_SpecB} is
state-dependent: its value shifts with the distributive
channel ($\omega_{t-1}$), the price channel
($p^\kappa_t$), and the empirical growth path
$d\kappa_t / dk^{NRC}_t$. Whether $\theta^{SpecB}_t$
exceeds or falls below unity is an empirical question
determined by the interaction of these channels under a
given accumulation regime. The identification strategy in
Section~3.4 specifies how this regime-dependence is
estimated and tested.



\subsection{The Peripheral Extension: Balance-of-Payments
	Constraint and the Corner Solution}

\paragraph{3.22}
In the capitalist periphery, the intensive margin of
accumulation is structurally tethered to the world market and
constrained by domestic social friction. Following
\citet{Prebisch1981} and \citet{Kaldor1959}, peripheral
capitalism is characterized by high consumption propensities
out of profits, structural inelasticities in agriculture, and
a chronic vulnerability to balance-of-payments crises
\citep{Ahumada2019}. This introduces a dual realization
bottleneck: the macroeconomic capacity to import and the
inelastic domestic supply of wage-goods. Section~2.5
established this constraint qualitatively as the peripheral
accumulation wedge. This subsection formalizes it as a strict
macroeconomic boundary constraint on the cost-minimizing
firm's choice of technique, integrating the structuralist
insights of \citet{Prebisch1981} into the heterogeneous
capital framework of Section~3.2.

\subsubsection{The Import-Biased Intensive Margin and the
	Balance of Payments Constraint}

\paragraph{3.23}
In Section~3.2, the cost-minimizing firm chooses the
mechanization gap $\kappa_t$ to maximize the net viability
surplus:
\begin{equation}
	\max_{\kappa_t}\; c(\kappa_t)
	= \Phi(\kappa_t) - \kappa_t \cdot MC(\omega_{t-1}, p^\kappa_t)
	\label{eq:33:objective_periphery}
\end{equation}
where $\Phi(\kappa_t) = \lambda_0 + \lambda_1 \kappa_t
+ \lambda_2 \kappa_t^2$ is the regime-conditioned
mechanization frontier, and $MC(\omega_{t-1}, p^\kappa_t)
= \gamma_0 - \gamma_1 \omega_{t-1} + \gamma_2 p^\kappa_t$
is the marginal viability cost of the gap. In the core, this
optimization proceeds unconstrained, yielding the optimal
mechanization gap $\kappa_t^*$. In the periphery, however,
active mechanization is strictly import-biased: widening the
mechanization gap requires imported capital goods, specialized
machinery, and technologically non-substitutable spare parts \citep{Vernengo2006}.

\paragraph{3.24}
\textbf{Assumption 3.} \textit{Import-biased mechanization.}
The physical import requirement per unit of output
($m^{req}_t$) is a linear function of the chosen technique:
\begin{equation}
	m^{req}_t = \nu \cdot \kappa_t
	\label{eq:33:import_req}
\end{equation}
where $\nu > 0$ captures the structural import-intensity of
the mechanization gap. Simultaneously, the physical capacity
to satisfy this requirement is bounded by the macroeconomic
\textbf{capacity to import} per unit of output
($m^{avail}_t = M^{avail}_t / Y_t$). Following
\citet{Prebisch1981}, peripheral export growth is governed by
the center's income elasticity of demand for primary goods
($\epsilon_{m,c} < 1$), structurally lagging behind domestic
import demand. This external asymmetry, compounded by high
consumption propensities out of capitalist profits
\citep{Kaldor1959}, severely squeezes the foreign exchange
available for productive accumulation. This imposes a strict
external ceiling on the choice of technique:
\begin{equation}
	\nu \cdot \kappa_t \le m^{avail}_t
	\implies \kappa_t \le \bar{\kappa}_t
	\equiv \frac{m^{avail}_t}{\nu}
	\label{eq:33:bop_ceiling}
\end{equation}

\subsubsection{The Macroeconomic Allocation of Foreign
	Exchange}

\paragraph{3.25}
\textbf{Assumption 4.} \textit{Balance-of-payments
	constraint.} Total import demand cannot exceed export
capacity. The total supply of foreign exchange per unit of
output is driven by primary exports
($x_t = p_t^x X_t / Y_t$). Due to the center's low income
elasticity for primary goods, $x_t$ is structurally rigid.
The total demand for foreign exchange per unit of output
consists of three competing components:
\begin{enumerate}
	\item \textbf{Wage-goods imports} ($m^W_t$): driven by the
	wage share $\omega_{t-1}$ and the structural inelasticity
	of domestic agriculture \citep{Kaldor1959}.
	$m^W_t = m_w \omega_{t-1}$.
	\item \textbf{Conspicuous consumption imports} ($m^C_t$):
	driven by the profit share $\pi_t$ and the high propensity
	of the peripheral upper strata to imitate the consumption
	patterns of the center \citep{Prebisch1981, Kaldor1959}.
	$m^C_t = m_c c_k \pi_t$.
	\item \textbf{Productive capital imports} ($m^K_t$): driven
	by the mechanization gap. $m^K_t = \nu \kappa_t$.
\end{enumerate}

\paragraph{3.26}
\textbf{Assumption 5.} \textit{Surplus leakage.} The
propensity to consume out of profits satisfies
$c_k > c_k^{\min}$, where $c_k^{\min}$ is the threshold
below which domestic savings suffice to finance capital
formation without external borrowing. Under the
developmentalist regime, the leakage channel is conspicuous
consumption imports by domestic elites
\citep{Prebisch1981, Kaldor1959}. Under the neoliberal
regime, the leakage channel is profit repatriation by
multinational corporations and capital flight
\citep{Ahumada2019}. In both cases, $c_k > c_k^{\min}$
ensures that the effective profit share available for
productive reinvestment is compressed below the nominal
level: $\pi_t^{eff} = \pi_t \cdot \phi_t < \pi_t$.

\paragraph{3.27}
The balance of payments constraint requires:
\begin{equation}
	m_w \omega_{t-1} + m_c c_k \pi_t + \nu \kappa_t \le x_t
	\label{eq:33:bop_full}
\end{equation}
Using the distributive identity $\pi_t = 1 - \omega_{t-1}$,
we substitute and solve for the maximum feasible mechanization
gap:
\begin{align}
	m_w \omega_{t-1} + m_c c_k (1 - \omega_{t-1})
	+ \nu \kappa_t &\le x_t
	\label{eq:33:bop_sub} \\
	\nu \kappa_t &\le x_t - m_c c_k
	+ (m_c c_k - m_w) \omega_{t-1}
	\label{eq:33:bop_solved} \\
	\bar{\kappa}_t(\omega_{t-1})
	&= \frac{1}{\nu} \left[ x_t - m_c c_k
	+ (m_c c_k - m_w) \omega_{t-1} \right]
	\label{eq:33:kappa_bar}
\end{align}
This equation formalizes the structuralist critique: the base
capacity to import capital goods is determined by export
earnings minus the baseline luxury consumption of the
capitalist class ($x_t - m_c c_k$). The sign of
$(m_c c_k - m_w)$ determines whether a rising wage share
expands or contracts the feasible mechanization gap.

\subsubsection{The Constrained Cost-Surplus Maximization
	Problem}

\paragraph{3.28}
The peripheral firm's tactical problem is therefore a
constrained cost-surplus maximization:
\begin{align}
	\max_{\kappa_t}\; c(\kappa_t)
	&= \Phi(\kappa_t) - \kappa_t \cdot MC(\omega_{t-1}, p^\kappa_t)
	\label{eq:33:constrained_obj} \\
	\text{subject to:}\quad & \kappa_t \le \bar{\kappa}_t
	\label{eq:33:constraint}
\end{align}
This formulation defines two distinct accumulation regimes,
endogenous to the interaction between domestic distributive
conflict and the global market.

\subsubsection{Endogenous Regime Shifts: The Corner Solution}

\paragraph{3.29}
The firm's unconstrained optimal technique
$\kappa_t^*(\omega_{t-1})$ is strictly increasing in the wage
share ($\partial \kappa_t^* / \partial \omega_{t-1} > 0$) due
to the Okishio labor-saving bias established in
\eqref{eq:32:cs_omega}.

\paragraph{3.30}
\textbf{Proposition 4.} \textit{Peripheral regime shifts.}
Two regimes are endogenous to the interaction between domestic
distributive conflict and the global market:
\begin{enumerate}
	\item \textbf{Regime A (Slack BoP):} When
	$\kappa_t^* \le \bar{\kappa}_t$, the firm solves the
	standard first-order condition, yielding $\kappa_t^*$. The
	labor-saving bias successfully extracts the Okishio
	distributive dividend.
	\item \textbf{Regime B (Peripheral realization trap):} When
	$\kappa_t^*(\omega_{t-1}) > \bar{\kappa}_t(\omega_{t-1})$,
	the constraint binds and the firm is forced into a corner
	solution:
	\begin{equation}
		\kappa_t = \bar{\kappa}_t
		\label{eq:33:corner}
	\end{equation}
\end{enumerate}
The corner solution arises through two intersecting structural
vulnerabilities. First, if $m_c c_k$ is large (the
conspicuous consumption drain), the baseline FX availability
is severely squeezed, lowering the intercept of
$\bar{\kappa}_t$. Second, if domestic agriculture is highly
inelastic ($m_w > m_c c_k$), a rising wage share
simultaneously increases the \textit{demand} for capital
imports (via $\kappa_t^*$) while \textit{contracting} the
supply of available FX (via $\bar{\kappa}_t$).
\textit{Proof.} Immediate from the Kuhn--Tucker conditions
applied to \eqref{eq:33:constrained_obj}--\eqref{eq:33:constraint}.
$\hfill\blacksquare$

\subsubsection{Effective Profit Compression and the Blocked
	Surplus}

\paragraph{3.31}
\textbf{Definition 3.} \textit{FX-Conversion Ratio.} The
FX-conversion ratio $\phi_t$ is the ratio of the maximum
feasible technique to the optimal unconstrained technique:
\begin{equation}
	\phi_t = \min\left(1, \frac{\bar{\kappa}_t}{\kappa_t^*}\right)
	\label{eq:33:phi}
\end{equation}
When the external constraint binds ($\kappa_t^* > \bar{\kappa}_t$),
$\phi_t < 1$. This generates a direct compression of the
effective profit share available for intensive accumulation.

\paragraph{3.32}
\textbf{Definition 4.} \textit{Blocked surplus.} Even if
domestic wage suppression yields a high nominal profit share
($\pi_t$), the physical inability to convert that surplus into
imported machinery squeezes the effective profit rate:
\begin{equation}
	\pi_t^{eff} = \pi_t \cdot \phi_t
	= (1 - \omega_{t-1}) \cdot \phi_t
	\label{eq:33:pi_eff}
\end{equation}
The difference constitutes the \textbf{blocked surplus}:
\begin{equation}
	\text{Blocked Surplus} = \pi_t(1 - \phi_t)
	\label{eq:33:blocked}
\end{equation}
Because the accumulation of machinery is driven by the
\textit{effective} profit share
($g_{K^{ME}} = \chi^{ME} \pi_t^{eff} - \delta^{ME}$), the
constraint decouples nominal profitability from capacity
formation. If capitalists squeeze wages to raise $\pi_t$, they
inadvertently increase conspicuous consumption ($m^C_t$),
which lowers $\bar{\kappa}_t$ and compresses $\phi_t$.

\subsubsection{The Regime-Shifting Aggregate
	Transformation Elasticity}

\paragraph{3.33}
\textbf{Definition 5.} \textit{BoP binding ratio.} The BoP
binding ratio $\tau_t$ is the ratio of productive capital
imports to import capacity:
\begin{equation}
	\tau_t \equiv \frac{\nu \cdot I^{ME}_t}
	{\text{ToT}_t \times X_t}
	\label{eq:33:tau}
\end{equation}
where $I^{ME}_t = K^{ME}_t - (1-\delta^{ME}) K^{ME}_{t-1}$
is gross investment in machinery and equipment, derived from
the perpetual inventory identity, and $\text{ToT}_t \times X_t$
proxies the import capacity (FX supply). Under Assumption~3,
productive capital imports $m^K_t = \nu \kappa_t$ are proxied
empirically by gross investment in machinery and equipment
$I^{ME}_t$, since mechanization is import-biased in the
periphery. The regime switch is governed by $\tau_t$:
\begin{itemize}
	\item $\tau_t < \bar{\tau}$: BoP constraint is
	\textbf{slack} (Regime A: unconstrained optimum).
	\item $\tau_t \ge \bar{\tau}$: BoP constraint is
	\textbf{binding} (Regime B: corner solution).
\end{itemize}
The threshold $\bar{\tau}$ is determined empirically by the
threshold regression in Section~4.

\paragraph{3.34}
The peripheral extension generates a regime-shifting
aggregate transformation elasticity. Under \textbf{Regime~A}
(slack BoP), the constraint does not bind, and the
elasticity coincides with the unconstrained expression from
Section~3.2:
\begin{equation}
	\theta^{A}_t = \theta^{SpecB}_t
	= \frac{\theta_{NRC}
		+ \left(\theta_\kappa + \theta_{\omega\kappa}\omega_{t-1}
		+ \theta_{\kappa p\kappa} p^\kappa_t\right)
		\frac{d\kappa_t^*}{dk^{NRC}_t}}
	{1 + s_{ME,t}\frac{d\kappa_t^*}{dk^{NRC}_t}}
	\label{eq:33:theta_A}
\end{equation}

\paragraph{3.35}
Under \textbf{Regime~B} (binding BoP), the firm is forced
into the corner solution $\kappa_t = \bar{\kappa}_t$. The
marginal capacity payoff at the corner is strictly positive
because the firm would like to mechanize further but cannot:
\begin{equation}
	MP_\kappa\big|_{\kappa=\bar{\kappa}}
	= \Phi'(\bar{\kappa}_t) - MC(\omega_{t-1}, p^\kappa_t) > 0
	\label{eq:33:MP_corner}
\end{equation}
The Kuhn--Tucker multiplier on \eqref{eq:33:constraint} is
strictly positive, confirming that the constraint binds. The
aggregate transformation elasticity under the binding regime
becomes:
\begin{equation}
	\theta^{B}_t
	= \frac{\theta_{NRC}
		+ MP_\kappa\big|_{\kappa=\bar{\kappa}}
		\cdot \frac{d\bar{\kappa}_t}{dk^{NRC}_t}}
	{1 + s_{ME,t}\frac{d\bar{\kappa}_t}{dk^{NRC}_t}}
	\label{eq:33:theta_B}
\end{equation}
Since $\bar{\kappa}_t$ is determined by the BoP constraint
\eqref{eq:33:kappa_bar} rather than by the first-order
condition \eqref{eq:32:foc_kappa}, the growth path
$d\bar{\kappa}_t / dk^{NRC}_t$ differs from the unconstrained
path $d\kappa_t^* / dk^{NRC}_t$. The FX-conversion ratio
$\phi_t = \bar{\kappa}_t / \kappa_t^* < 1$ compresses the
intensive margin, and the elasticity under the binding regime
is strictly lower:
\begin{equation}
	\theta^{B}_t < \theta^{A}_t
	\label{eq:33:theta_inequality}
\end{equation}

\paragraph{3.36}
\textbf{Proposition 5.} \textit{Regime-shifting elasticity.}
The peripheral aggregate transformation elasticity is a
regime-dependent parameter:
\begin{equation}
	\theta^{periphery}_t
	= \theta^{A}_t \cdot \mathbf{1}\{\kappa_t^* \le \bar{\kappa}_t\}
	+ \theta^{B}_t \cdot \mathbf{1}\{\kappa_t^* > \bar{\kappa}_t\}
	\label{eq:33:theta_regime}
\end{equation}
The elasticity shifts discretely when the BoP constraint
becomes binding. The regime switch is governed by the
interaction between the desired mechanization gap
$\kappa_t^*(\omega_{t-1}, p^\kappa_t)$ and the maximum
feasible gap $\bar{\kappa}_t(\omega_{t-1})$ determined by
export capacity, conspicuous consumption, and wage-goods
imports. An increase in the wage share $\omega_{t-1}$
simultaneously raises $\kappa_t^*$ (via the Okishio
labor-saving bias) and can contract $\bar{\kappa}_t$ (via
the wage-goods import channel when $m_w > m_c c_k$), making
the regime switch more likely. This is the peripheral
realization trap: distributive pressure that would induce
productive mechanization in the center instead triggers the
BoP constraint in the periphery, truncating the
transformation elasticity.

\subsubsection{The Failure of the Okishio Mechanism and
	Structural Inflation}

\paragraph{3.37}
\textbf{Proposition 6.} \textit{Blocked Okishio mechanism.}
At the boundary $\kappa_t = \bar{\kappa}_t$, the
implementation of the Okishio viability condition is blocked.
The firm satisfies the viability criterion---the technique
would reduce unit costs---but cannot deploy it because the
required machinery cannot be imported. The wage-savings
dividend cannot be capitalized into physical capacity. The
firm suffers a persistent shortfall in its cost-surplus
relative to its unconstrained optimal state.
\textit{Proof.} At $\kappa_t = \bar{\kappa}_t$, the
Kuhn--Tucker multiplier on \eqref{eq:33:constraint} is
strictly positive, so the first-order condition
$\partial c / \partial \kappa_t = 0$ does not hold. The
marginal viability yield exceeds the marginal cost at the
boundary, but the firm cannot expand further.
$\hfill\blacksquare$

\paragraph{3.38}
Unable to extract the distributive dividend via mechanization,
capitalists attempt to recover the blocked surplus through the
only remaining channel: price markups. When the wage-goods
bottleneck forces nominal wage adjustments and the external BoP
bottleneck blocks defensive mechanization, any further increase
in the wage share directly crushes profitability. The resulting
defensive pricing behavior ignites the structural inflation
dynamic \citep{Sunkel1958, Prebisch1981}. This is the
peripheral realization trap: nominal profitability coexists
with truncated mechanization and stalled structural envelopes,
permanently destabilizing the long-run coupling between plant
and machinery.


\subsection{Synthesis: What the Three Models Jointly
	Demonstrate}

\paragraph{3.38}
The three models developed in this section formalize a single
theoretical claim: that decentralized, individually
cost-minimizing decisions generate aggregate
disproportionality under the unbalanced growth closure
\citep{Okishio1961, Basu2026}.
Table~\ref{tab:coherence} summarizes the articulation
between the Okishio and Vidal frameworks that underpins
this claim.

\begin{table}[H]
	\centering
	\small
	\caption{Theoretical Coherence: Market Anarchy and the Internally Divided Firm Articulation Across Three Models}
	\label{tab:coherence}
	\begin{tabular}{p{2.2cm}p{3.8cm}p{3.8cm}p{4.2cm}}
		\toprule
		\textbf{Element} & \textbf{Descentralized Accumulation} &
		\textbf{Internally Divided Firm} & \textbf{Formal model captures} \\
		\midrule
		\midrule
		Micro logic & Viability condition (cost
		reduction at prevailing prices) & Internal
		contradiction (control vs.\ cooperation) &
		Optimization problem = cost-minimizing pole;
		constraints = cooperation limits \\
		Macro logic & Decentralized decisions $\to$
		disproportionality & Divided firm $\to$
		contested operating conditions & $\theta \neq 1$
		captures gap between accumulation and capacity \\
		Disequilibrium & Unbalanced growth (supply
		$\neq$ demand) & ``Mediocre sufficiency''
		\citep{Vidal2022} (no coherent optimum) & $\mu$
		demonstrates gap; $\theta$ captures structural
		inefficiency \\
		Peripheral specificity & BoP constraint blocks
		\textit{implementation} of viability & External
		constraint compounds internal division & Corner
		solution: $\kappa = \bar{\kappa}$ when
		constraint binds \\
		\bottomrule
	\end{tabular}
\end{table}

\paragraph{3.39}
The formal model does not assert equilibrium. It
demonstrates the negative: that capitalist accumulation,
left to its own decentralized logic, does not generate
proportional reproduction. The capacity utilization rate
$\mu$ is the empirical signature of this
disproportionality: it shows that the normal operating
conditions of capitalist firms, determined by their
cost-minimizing mechanization choices, do not meet
aggregate demand conditions. The transformation elasticity
$\theta$ captures this structural gap between accumulation
and realized capacity. Whether $\theta$ exceeds or falls
below unity is an empirical question determined by the
interaction of distributive and price channels under a
given accumulation regime. In the center, this gap is
cyclical---mediated by institutional configurations that
can be exhausted and reconstituted
\citep{Vidal2015, Vidal2019}. In the periphery, this gap
is permanent---locked in by the structural ceiling of
external subordination \citep{Prebisch1981, Ahumada2019}.

\paragraph{3.40}
Section~4 consolidates the empirical identification
strategy. It translates the three analytical models into
estimable cointegrating specifications, derives the
threshold regression for the peripheral corner solution,
and specifies the capacity recovery protocol. The
empirical analysis bounds this multi-sector mechanism to
a single-sector aggregate, which masks internal sectoral
divergences but preserves the core transformation logic.
\FloatBarrier
\section{Data, Measurement, and Empirical Results}
\subsection{Data and Measurement}
\label{sec:41:data_measurement}

% -------------------------------------------------
% 4.1.1 THE DUAL VALUATION DISTINCTION
% -------------------------------------------------
\subsubsection{The Dual Valuation Distinction}
\label{sec:411:dual_valuation}

This section constructs the capital stock series for the
choice-of-technique models of \cref{sec:3:analytical_framework}.
The BEA Fixed Assets Accounts report a single current-cost net stock.
That series treats the physical capacity of installed capital and
the monetary value of capital advanced as one magnitude, and the
conflation introduces a spurious correlation between capital
intensity and the wage share: as mechanization raises capital
intensity, Consumption of Fixed Capital (CFC) rises in the GVA
denominator, depressing GVA-basis wage shares even when the rate of
surplus-value extraction is unchanged. To eliminate this bias, I
reconstruct the stock account-by-account using the Generalized
Perpetual Inventory Method (GPIM) of \citet[App.~6.5]{Shaikh2016},
which builds capacity and profitability measures from direct nominal
investment flows and asset-specific price indexes.

\paragraph{The dual-valuation distinction.}
Estimates of capacity utilization under my unbalanced growth framework, require distinct 
capital stock variables. The capital value initially invested in plant and equipment returns
gradually to its money form as fixed assets depreciate; the
accumulated depreciation allowances remain part of total capital
advanced regardless of their monetary form, because their recovery
is what permits the reproduction of the capitalist enterprise 
\citep[App.~6.3]{Shaikh2016}. I therefore distinct two different types of
measurements of capital stocks:

\begin{enumerate}
	\item \textbf{Capital stock for capacity estimation: gross real
		GPIM stock.}
	Productive capacity depends on capital goods still physically in
	operation. The capital stock measure adequate for estimating
	capacity is the gross real stock $K^{G,R}_{j,t}$, weighted
	exclusively by a physical survival schedule. Productive capacity
	$Y^p_t$ is itself a latent variable, recovered through the
	estimation of the transformation elasticity $\theta$
	(\cref{sec:3:analytical_framework}); the gross real GPIM stock is
	the accumulation variable $k_t$ included in my estimation
	equation, not a measure of capacity itself. A machine whose book value has
	fallen to zero, but remains in service until it is physically retired, still should be accounted as physical capital in operation.  The OECD manual confirms that the gross stock, valued at ``as
	new'' prices, ``is widely used as a broad indicator of the
	productive capacity of a country'' \citep{OECD2001}.
	
	\item \textbf{Profitability measure: net current-cost stock.}
	Profitability is a monetary ratio between a profit flow and the
	value of capital advanced. The relevant measure is the net
	current-cost stock $K^{N,V}_{j,t}$, which tracks the monetary
	value of capital after financial depreciation and write-downs.
	This measure serves as the anchor for the estimation of an implicit 
	price of capital goods under a stock-flow-consistent framework but is not included within the structural estimation of non-observed productive capacities. 
\end{enumerate}

A single aggregate net stock serving both purposes is common in
empirical macroeconomics. This practice introduces a mechanical
correlation between the capital stock and the depreciation flow,
introducing potential biases to identify distributive effects on productive capacities formation. 

\paragraph{Valuation choice within the GPIM.}
The dual-valuation distinction does not require imposing a single
valuation convention throughout. \citet[App.~6.3]{Shaikh2016}
advocates the gross stock as the denominator for both the profit
rate and the capacity ratio. The GPIM he develops, however,
constructs the survival schedule $S_j(v)$, the age-price schedule
$(1-d_j)^v$, and the SFC implicit price index $P_{j,t}$ as distinct
components within a single stock-flow-consistent recursion. This
structure permits the researcher to select the valuation for each
analytical purpose without abandoning accounting coherence.

For the profit rate, the denominator must measure the monetary value
of capital advanced. The net current-cost stock serves this role: it
tracks the book value of capital that the firm must recover to
continue production. A machine that has depreciated to a low book
value may still operate at full physical capacity, but the firm's
cost-recovery problem is defined by the value it must recoup, not by
the physical state of the equipment \citep{Basu2013}. For estimates
related to capacity utilization, the capital stock input must reflect
the physical capital still in operation. The gross real GPIM stock,
weighted exclusively by the survival schedule $S_j(v)$ and excluding
the age-price factor $(1-d_j)^v$, serves this role.

The quantitative significance of the gross/net distinction is not
confined to the center. \citet{Ffrench2016} show that in Chile
depreciation absorbs between 8 and 13 percent of GDP and that net
investment is barely half of gross investment in the most recent
period. Service lives of machinery have declined from 20 years in
the 1940s to 11 years by 2005, reflecting accelerating obsolescence.
They conclude that the capital stock entering a production function
must be the gross stock of physically surviving assets, not the
depreciated monetary residual.

% -------------------------------------------------
% 4.1.2 U.S. DATA PROVENANCE AND CAPITAL RECONSTRUCTION
% -------------------------------------------------
\subsubsection{Gross Capital Stock Reconstruction and Measurement in the United States}
\label{sec:412:us_capital}

\paragraph{The Generalized Perpetual Inventory Method.}
In this part I present the procedures of construction for 
the capital stock series for the United States. Machinery (ME) and nonresidential structures (NRC) do not
share a natural physical unit; they become comparable only through a
monetary-mediated accounting procedure. I reconstruct the stock
account-by-account rather than aggregating nominal values and
deflating with a generic index. The GPIM accumulation rules for
historical-cost, current-cost, and constant-dollar stocks
\citep[App.~6.5]{Shaikh2016} are derived in detail in
Appendix~\ref{app:A1:gpim_protocol}; here I state the construction.

Physical survival of capital vintages follows a Weibull retirement
distribution. For asset class $j \in \{\text{ME}, \text{NRC}\}$,
the proportion of a vintage surviving after $v$ years is:
\begin{equation}
	S_j(v) = \exp\!\left[ -\left( \frac{v}{\lambda_j}
	\right)^{\alpha_j} \right]
	\label{eq:41:weibull_survival}
\end{equation}
where $\alpha_j$ is the Weibull shape parameter and the scale
parameter is $\lambda_j = L_j / \Gamma(1 + 1/\alpha_j)$, derived
from the mean service life $L_j$. Unlike geometric decay models that
continuously erode the physical stock, the Weibull curve ensures
that capital goods remain fully operative until they are physically
removed from service. Empirical studies of directly observed asset
discards confirm that the Weibull family provides a flexible
approximation of retirement patterns. \citet{Nomura2005}, using
Japanese micro-data on 66 asset classes, finds that almost half of
assets exhibit a progressively increasing hazard rate
($\alpha > 1$). The calibrated parameters $L_{\text{ME}} = 14$,
$\alpha_{\text{ME}} = 1.7$ and $L_{\text{NRC}} = 30$,
$\alpha_{\text{NRC}} = 1.6$ follow the BEA Fixed Assets methodology
\citep{Fraumeni1997}.

Physical survival alone cannot identify the SFC implicit price
index. The BEA current-cost net stock anchors reflect both physical
retirement and age-price decay. To recover the asset-specific price
index, I impose an externally documented geometric declining-balance
age-price profile:\footnote{The BEA's current methodology employs
	geometric depreciation, which assumes that the efficiency pattern and
	the depreciation pattern have the same form \citep{Fraumeni1997}.
	The GPIM construction separates these two dimensions: physical
	survival follows a Weibull retirement distribution, while financial
	age-price decay follows a geometric profile. This separation is
	necessary to construct distinct capacity and profitability measures.}
\begin{equation}
	V_j(v) = S_j(v) \cdot (1 - d_j)^v
	\label{eq:41:net_value_schedule}
\end{equation}
where $d_j$ is the asset-specific age-price decay rate. The
declining-balance formula applied to the calibrated service lives
yields $d_{\text{ME}} = 0.110$ and $d_{\text{NRC}} = 0.024$,
consistent with the geometric rates documented by the BEA
\citep{Fraumeni1997}.

The SFC implicit price index $P_{j,t}$ is the price level that
closes the stock-flow identity for the net current-cost stock. In
each period, the BEA anchor must satisfy
\begin{equation}
	K^{\text{NET,CC}}_{j,t}
	= I^{\text{NOM}}_{j,t}
	+ \frac{P_{j,t}}{100}
	\sum_{v=1}^{L_j} V_j(v)\, I^{\text{REAL}}_{j,t-v}
	\label{eq:41:sfc_identity}
\end{equation}
where $I^{\text{REAL}}_{j,t-v} = I^{\text{NOM}}_{j,t-v}/(P_{j,t-v}/100)$
is real investment deflated by the previously recovered SFC price.
The first term on the right-hand side is current-year nominal
investment; the second is the net-value-weighted sum of surviving
real vintages revalued at the current price. The price index is
``implicit'' because it is not directly observed: it is the unique
value that, applied to the survival-weighted real vintages,
reproduces the BEA's reported current-cost net stock. Solving
\eqref{eq:41:sfc_identity} for $P_{j,t}$:
\begin{equation}
	P_{j,t} = 100 \cdot
	\frac{K^{\text{NET,CC}}_{j,t} - I^{\text{NOM}}_{j,t}}
	{\displaystyle\sum_{v=1}^{L_j} V_j(v)\, I^{\text{REAL}}_{j,t-v}}
	\label{eq:41:sfc_price_recursion}
\end{equation}
The forward recursion proceeds period by period: given
$P_{j,t-1}, P_{j,t-2}, \ldots$, the current price $P_{j,t}$ is
uniquely determined by the requirement that the constructed stock
reproduces the BEA anchor.
Appendix~\ref{app:A1:gpim_protocol} documents the recursion, the
initial conditions, and the base-year normalization. The BEA's Fisher
quantity index is retained solely as a validation comparison and does
not enter the GPIM baseline.

The capital stock used to estimate productive capacity is the gross
real GPIM stock, constructed as the survival-weighted sum of real
investment vintages:
\begin{equation}
	K^{G,R}_{j,t} = \sum_{v=0}^{L_j} S_j(v)\,
	I^{\text{REAL}}_{j,t-v}
	\label{eq:41:gross_gpim_stock}
\end{equation}
where $I^{\text{REAL}}_{j,t-v}$ is real investment deflated by the
previously recovered SFC price. Because
\cref{eq:41:gross_gpim_stock} relies exclusively on the survival
schedule $S_j(v)$ and excludes the age-price factor $(1-d_j)^v$, it
measures the constant-price purchasing power embodied in
heterogeneous capital goods still in physical operation. Financial
write-downs, revaluation gains, and depreciation allowances do not
enter.

\paragraph{Validation and sample scope.}
The GPIM construction closes the stock-flow identity exactly: the
maximum absolute reconstruction residual across both asset classes
is $10^{-6}$ million dollars, confirming that the recovered SFC
implicit price indexes reproduce the BEA current-cost net stock
anchors. All GPIM series are normalized to $2017 = 100$.
The capital stock series has a 30-year warm-up period (1901--1930).
The GPIM stock at any year $t$ requires investment vintages
extending back $L_j$ years. Because the nominal investment series
begins in 1901, the vintage summation for the longest-lived asset
class (nonresidential structures, $L_{\text{NRC}} = 30$) is
truncated for any year before 1931. The stock values for
1901--1930 are therefore incomplete, omitting pre-1901 vintages
that would still be physically in service. I restrict the
estimation sample to 1931--2024, where the full 30-year vintage
window is populated, yielding $T = 94$ annual observations.
The sectoral scope is restricted to nonfinancial corporate (NFC)
business, excluding non-profit institutions serving households
(NPISH), government, and owner-occupied housing. This sectoral
restriction isolates the private productive business sector from
public infrastructure and imputed residential rents.

\subsubsection{Distributive Conflicts and Wage Share Dynamics in the United States}


\subsubsection{Capital Series Harmonization and Historical Sources for Chile}
\label{sec:413:chile_sfc}

This subsection describes the stock-flow consistent (SFC) methodology
used to reconstruct gross and net capital stocks for Chile,
1900--1994, in constant 2003 Chilean pesos (CLP). The procedure
enforces the accumulation identities for gross and net stocks
jointly, yields implicit depreciation and retirement rates for each
asset class, and delivers coherent series for machinery and equipment
(ME), nonresidential construction (NRC), residential construction
(RC), and their aggregates. The construction proceeds from a single
investment control total, disaggregated by asset, accumulated via
Weibull survival functions, and linked to net stocks through a
gap-ratio approach anchored to the \citet{Hofman2000b} overlap
window.

\paragraph{The stock-flow identities.}
The reconstruction is governed by two laws of motion. The gross
stock evolves according to
\begin{equation}
	\Delta K^{g}_{j,t} = I_{j,t} - R_{j,t}
	\label{eq:42:gross_identity}
\end{equation}
where $I_{j,t}$ is gross fixed investment and $R_{j,t}$ denotes
physical retirements (scrapping of obsolete vintages, demolition,
and technological withdrawal). The net stock evolves according to
\begin{equation}
	\Delta K^{n}_{j,t} = I_{j,t} - D_{j,t}
	\label{eq:42:net_identity}
\end{equation}
where $D_{j,t}$ is depreciation (the monetary consumption of fixed
capital). Combining \eqref{eq:42:gross_identity} and
\eqref{eq:42:net_identity} yields the joint SFC identity
\begin{equation}
	\Delta K^{g}_{j,t} + R_{j,t}
	= \Delta K^{n}_{j,t} + D_{j,t}
	\label{eq:42:joint_identity}
\end{equation}
which constrains the combination of depreciation and retirements at
every $t$ and for every asset $j$. Given observed stocks and
investment, the implied flows are recovered as $D_{j,t} = I_{j,t}
- \Delta K^{n}_{j,t}$ and $R_{j,t} = I_{j,t} - \Delta K^{g}_{j,t}$.
The capital gap $B_{j,t} \equiv K^{g}_{j,t} - K^{n}_{j,t}$ measures
the cumulative difference between installed capacity and book value.

\paragraph{The investment control total and the lambda bridge.}
The aggregate investment control total is drawn from the
\citet{PerezEyzaguirre2019} reconstruction of gross fixed capital
formation (GFCF), expressed in constant 2003 CLP. The control
concept is machinery plus construction, explicitly excluding
inventory variation.\footnote{Inventories are excluded from
	fixed-capital accumulation because the reconstruction models fixed
	capital only. Inventory investment does not enter the perpetual
	inventory recursion.}

Pérez-Eyzaguirre report investment in a price space and asset
aggregation that differ from the \citet{Hofman2000a} benchmark. The
two sources are linked by a level-normalization factor $\lambda$,
which corrects for definitional differences between the two series
rather than converting between price bases. The factor is estimated
as the median of the annual ratio of Hofman to Pérez investment over
the 95-year estimation window:
\begin{equation}
	\lambda = \mathrm{median}\!\left(
	\frac{I^{\text{Hofman}}_t}{I^{\text{Pérez}}_t}
	\right)_{t=1900}^{1994}
	= 4.1054
	\label{eq:42:lambda}
\end{equation}
The median rule is chosen for robustness to outliers in the early
decades, where source quality is weaker. The scaled aggregate is
then disaggregated into ME, NRC, and RC using normalized
asset-specific investment shares from \citet{Hofman2000a}:
\begin{equation}
	I_{j,t} = s_{j,t} \cdot \lambda \cdot I^{\text{Pérez}}_t,
	\quad
	\sum_{j \in \{\text{ME},\, \text{NRC},\, \text{RC}\}}
	s_{j,t} = 1
	\label{eq:42:disaggregation}
\end{equation}
for all $t$.

\paragraph{Weibull survival and gross stock accumulation.}
Physical survival of capital vintages follows a Weibull retirement
distribution. For asset class $j \in \{\text{ME}, \text{NRC},
\text{RC}\}$, the proportion of a vintage surviving after $v$ years
is
\begin{equation}
	S_j(v) = \exp\!\left[
	-\left( \frac{v}{\lambda_j} \right)^{\alpha_j}
	\right]
	\label{eq:42:weibull_chile}
\end{equation}
where $\alpha_j$ is the Weibull shape parameter and the scale
parameter is $\lambda_j = L_j / \Gamma(1 + 1/\alpha_j)$, derived
from the mean service life $L_j$. The survival parameters are
estimated from the \citet{Hofman2000b} overlap window
(1950--1994), in which both gross and net stocks are observed, by
fitting Weibull curves to the implied retirement patterns recovered
from the stock-flow identities.
Table~\ref{tab:42:survival_params} reports the calibrated
parameters.

\begin{table}[h]
	\centering
	\caption{Weibull survival parameters, Chile (1900--1994)}
	\label{tab:42:survival_params}
	\begin{tabular}{lccc}
		\toprule
		Asset & Scale ($\lambda_j$) & Shape ($\alpha_j$) & Mean life
		($L_j$, years) \\
		\midrule
		ME  & 15.90 & 2.42 & 14.1 \\
		NRC & 41.30 & 5.00 & 37.9 \\
		RC  & 51.16 & 5.00 & 47.0 \\
		\bottomrule
	\end{tabular}
\end{table}

All shape parameters exceed unity, indicating progressively
increasing hazard rates: capital goods remain fully operative in
early years and are withdrawn from service at an accelerating rate
as they approach the mean service life. The shape parameters for
NRC and RC are constrained at 5.0, reflecting the long and
relatively uniform service lives of structures. The gross stock is
accumulated as the survival-weighted sum of real investment
vintages:
\begin{equation}
	K^{g}_{j,t} = \sum_{v=0}^{L_j} S_j(v)\, I_{j,t-v}
	\label{eq:42:gross_stock_chile}
\end{equation}
Pre-1950 stocks are anchored to the observed Hofman levels at 1950
via a multiplicative factor that matches $K^{g}_{j,1950}$ by asset.

\paragraph{The gap-ratio approach and net stock construction.}
Unlike the U.S. GPIM, which recovers the net stock via a separate
geometric age-price decay rate and an SFC implicit price recursion,
the Chile reconstruction derives the net stock from the observed
gross-net relationship in the \citet{Hofman2000b} overlap window.
The capital gap ratio
\begin{equation}
	\gamma_j \equiv \frac{B_j}{K^{g}_j}
	= \frac{K^{g}_j - K^{n}_j}{K^{g}_j}
	\label{eq:42:gap_ratio}
\end{equation}
is computed by asset from the \citet{Hofman2000b} data over the
overlap period (1950--1994). Two regimes are distinguished: a
``pre'' ratio (median of years $\leq 1954$) and a ``post'' ratio
(median of years $\geq 1990$), allowing for structural shifts in
the gross-net relationship over the century. Outside the overlap
window, the gap ratio is extended at the regime-specific median.
The net stock is then recovered as
\begin{equation}
	K^{n}_{j,t} = K^{g}_{j,t} \cdot (1 - \gamma_j)
	\label{eq:42:net_stock_chile}
\end{equation}
This approach inherits the gross-net relationship observed in the
Hofman accounting reference rather than imposing an external
depreciation convention. It is appropriate for Chile because no
independent age-price decay data are available for the historical
period; the \citet{Hofman2000b} overlap provides the only direct
evidence on the joint behavior of gross and net stocks.

\paragraph{Implicit depreciation and retirement rates.}
Given the reconstructed stocks and investment flows, the implicit
annual rates are defined as
\begin{equation}
	z_{j,t} = \frac{R_{j,t}}{K^{g}_{j,t-1}},
	\qquad
	\delta_{j,t} = \frac{D_{j,t}}{K^{n}_{j,t-1}}
	\label{eq:42:implicit_rates}
\end{equation}
where $z_{j,t}$ is the retirement (depletion) rate and
$\delta_{j,t}$ is the depreciation rate. These rates are not
imposed as parameters; they are derived from the data via the
stock-flow identities
\eqref{eq:42:gross_identity}--\eqref{eq:42:net_identity}. They
serve as diagnostic indicators of capital durability and
replacement dynamics, and are used in
\cref{sec:42:us_aggregate} to validate the internal consistency of
the reconstruction.

\paragraph{Validation and estimation sample.}
The SFC construction is validated through four checks. First,
additivity: all composite stocks and flows satisfy
$K_{\text{NR}} = K_{\text{ME}} + K_{\text{NRC}}$ and
$K_{\text{T}} = K_{\text{ME}} + K_{\text{NRC}} + K_{\text{RC}}$
at every $t$. Second, non-negativity: $K^{g}_{j,t} \geq 0$,
$K^{n}_{j,t} \geq 0$, and $B_{j,t} \geq 0$ for all assets and
years. Third, rate bounds: $z_{j,t} \in [0,1]$ and
$\delta_{j,t} \in [0,1]$ for all observations from 1950 onward.
Fourth, identity residuals: $\Delta K^{g}_{j,t} - (I_{j,t} -
R_{j,t}) = 0$ and $\Delta K^{n}_{j,t} - (I_{j,t} - D_{j,t}) = 0$
to numerical precision. The estimation sample spans 1900--1994,
yielding $T = 95$ annual observations. The overlap window
(1950--1994) provides 45 years of direct \citet{Hofman2000b}
observations against which the reconstructed series are audited.

\paragraph{The BCCh forward extension (1995--2024).}
Extending the reconstruction beyond 1994 requires integrating the
official national accounts published by the Central Bank of Chile
(BCCh). A direct chain-linking of the BCCh series to the historical
baseline would introduce a level discontinuity, because the two
sources differ in base year, valuation conventions, and asset
classification. To eliminate this discontinuity while preserving the
validated 1900--1994 history, I employ a multiplicative splicing
procedure that fixes the legacy panel in place.

For the overlap year 1994, an asset-specific bridge factor is
computed as the ratio of the historical net stock to the BCCh
reported net stock:
\begin{equation}
	\beta_j = \frac{K^{n,\text{hist}}_{j,1994}}
	{K^{n,\text{BCCh}}_{j,1994}},
	\quad j \in \{\text{ME},\, \text{NRC},\, \text{RC}\}
	\label{eq:42:bridge_factor}
\end{equation}
The same factor $\beta_j$ is applied to both the BCCh net stock and
the BCCh depreciation series for all years 1995--2024. This
preserves the depreciation rates and asset dynamics implied by the
BCCh while maintaining continuity with the constant 2003 CLP
historical baseline. Table~\ref{tab:42:bridge_factors} reports the
estimated bridge factors.

\begin{table}[h]
	\centering
	\caption{Asset-specific bridge factors, 1994}
	\label{tab:42:bridge_factors}
	\begin{tabular}{lcc}
		\toprule
		Asset & Bridge factor ($\beta_j$) & Interpretation \\
		\midrule
		ME  & 5.888 & BCCh understates historical ME stock \\
		NRC & 1.033 & Near parity \\
		RC  & 1.050 & Near parity \\
		\bottomrule
	\end{tabular}
\end{table}

The large bridge factor for machinery and equipment reflects the
substantially lower level at which the BCCh reports ME stocks
relative to the Pérez-based historical reconstruction, a divergence
driven by differences in asset coverage and valuation methodology.
For nonresidential and residential construction, the bridge factors
are close to unity, indicating that the two sources agree
substantially on the level of structures.

Post-1994 real investment is inferred via the stock-flow identity:
\begin{equation}
	I_{j,t} = \Delta K^{n}_{j,t} + D_{j,t},
	\quad t \geq 1995
	\label{eq:42:post94_investment}
\end{equation}
where $D_{j,t}$ is the bridge-adjusted BCCh depreciation. The gross
stock and physical retirements for the extended period are then
reconstructed by applying the survival parameters from
Table~\ref{tab:42:survival_params} to the complete 1900--2024
investment series via
\cref{eq:42:gross_stock_chile}. This ensures that the gross stock
remains internally consistent across the splice boundary.


\subsubsection{Distributive Conflicts and Wage Share Dynamics in Chile}

\label{sec:41x:chile_distrib}
The distributive regressor in the transformation elasticity is the
wage share, $\omega_t$. For Chile, the wage share is drawn from the
revised dataset of real wages in Latin America constructed by
\citet{Astorga2017} and updated by \citet{Astorga2023}. The dataset
provides yearly series of real wages by skill level---unskilled
($w_4$), semi-skilled ($w_3$), and skilled ($w_2$)---together with
an economically-active-population (EAP) weighted average ($w$),
expressed in monthly purchasing-power-parity dollars at 1970
prices.\footnote{The revised dataset is publicly available as
	supplementary material to \citet{Astorga2023}.}

The wage share is computed as the wage-to-income ratio
$\omega_t = w_t / y_t$, where $y_t$ is real income per worker
derived from national income estimates. The functional distribution
panel extends the \citet{Astorga2023} coverage to 2024 by splicing
the revised wage series with post-2011 observations from official
Chilean statistical sources. The resulting $\omega_t$ series spans
1920--2024, yielding $T = 105$ annual observations.

\subsubsection{Aggregate Demand, Price Indices, and External Constraint variables for Chile}
\label{sec:415:chile_demand}

The aggregate demand variables---real GDP ($Y_t$), exports FOB
($X^{\text{FOB}}_t$), and gross machinery investment
($I_{\text{ME},t}$)---are drawn from the aggregate demand
reconstruction of \citet{PerezEyzaguirre2019}, expressed in constant
2003 Chilean pesos. This series revises and extends the historical
national accounts backbone of La República en Cifras
\citep{LudersDiazWagner2016}, correcting additivity and valuation
problems in the original Clio-Lab series while preserving the
sectoral GDP decomposition. Coverage spans 1860--2024.

The price and trade indices used to construct the relative capital
price proxies ($p^{\kappa}_t$), the capacity to
import ($M^{\text{cap}}_t$), are drawn from \textit{La República en Cifras}
\citep{LudersDiazWagner2016}: the real exchange rate index, the US wholesale capital price index, and the
terms of trade index. All three series are rebased to
2003 $= 100$ for internal consistency with the capital stock
series.\footnote{The terms of trade index extends back to 1810;
	the real exchange rate and US wholesale price index cover
	1900--2024.}

The capital import demand imbalance, combines sources from \textit{La República en Cifras} with my data set of harmonized capital stocks through a log difference of $log({I_{ME}}_t)$ - $log(M^{\text{cap}}_t)$, which reconciles the supply export-side identified by the index of imports capacities $(M^{\text{cap}}_t)$ of the external constraint, with the internal demand for foreign capital goods providing a both a more accurate variable for the regime shifting empirical strategy detailed below, while also a variant that allows me to check for roboustness of my estimates. 

A complementary data source for external sector variables are drawn from the World Historical Balance of Payments Database (WBOP) developed by
\citet{NievasPiketty2025}. The WBOP data set provides provides annual series on global trade flows and the world balance of payments---including goods,
services, income, and transfers---covering 57 countries over
the 1800--2025 period. Three series are used for Chile: the ratio
of manufactured goods imports to GDP (a proxy for capital goods
import dependence), net foreign income outflows as a share of GDP (a
proxy for profit remittance drain), and the net current account
balance as a share of GDP (an external balance indicator). All three
series are expressed in current US dollars at market exchange rates
and are available at the WBOP online repository.\footnote{All WBOP
	series are available at \texttt{wbop.world} and are also hosted in
	the World Inequality Database (\texttt{wid.world}), together with a
	detailed replication package including raw data sources, methods,
	and codes.}




%%%%%%%%%%%%%%%%%%%%%%%%%%%%%%%%%%%%%%%%%%%%%%%%%%%%%%%
%%%%% 4.2. Econometric Strategy IM-OLS Cointegration
%%%%%%%%%%%%%%%%%%%%%%%%%%%%%%%%%%%%%%%%%%%%%%%%%%%%%%%

\subsection{Econometric Strategy: Integrated Modified OLS and Cointegration}



%%%%%%%%%%%%%%%%%%%%%%%%%%%%%%%%%%%%%%%%%%%%%%%%%
%%%%% 4.3. US Results Aggregate Capital
%%%%%%%%%%%%%%%%%%%%%%%%%%%%%%%%%%%%%%%%%%%%%%%%%


\subsection{United States Empirical Evidence: Aggregate Capital Accumulation and Capacity Utilization}

\subsubsection{Baseline Empirical Estimates}

\subsubsection{Transformation Elasticity and Capacity Utilization Estimates}


%%%%%%%%%%%%%%%%%%%%%%%%%%%%%%%%%%%%%%%%%%%%%%%%%
%%%%% 4.4. US Results Aggregate Capital
%%%%%%%%%%%%%%%%%%%%%%%%%%%%%%%%%%%%%%%%%%%%%%%%%


\subsection{United States Empirical Evidence: Capital Heterogeneity and the Mechanization Gap}


\subsubsection{Heterogeneous Capital Stocks Empirical Estimates}


\subsubsection{Aggregate Transformation Elasticity and Capacity Utilization Estimates with Heterogeneous Capital Stocks}


%%%%%%%%%%%%%%%%%%%%%%%%%%%%%%%%%%%%%%%%%%%%%%%%%
%%%%% 4.5. Chile Results 
%%%%%%%%%%%%%%%%%%%%%%%%%%%%%%%%%%%%%%%%%%%%%%%%%
\subsection{Chilean Empirical Evidence: The Peripheral and External Constraints}

\subsubsection{Econometric Strategy: Threshold Conditional Mixed Process Estimation}


\subsubsection{Empirical Estimates of the Balance-of-Payments Constraint}

\subsubsection{Transformation and Capacity Utilization und Financial Subordination}



%%%%%%%%%%%%%%%%%%%%%%%%%%%%%%%%%%%%%%%%%%%%%%%%%
%%%%% 4.6 Comparative and Historical Analysis 
%%%%%%%%%%%%%%%%%%%%%%%%%%%%%%%%%%%%%%%%%%%%%%%%%
\subsection{Comparative Analysis and Theoretical Synthesis}



\FloatBarrier
\section{Comparative Determination: Center and Periphery}

\paragraph{A}
The empirical estimations established in Section 4 demonstrate that capacity utilization cannot be understood as a universal engineering constant, an unmediated behavioral ratio, or an empirical residual. Instead, utilization is a structurally reconstructed object that emerges from the identification of the aggregate transformation elasticity $\theta_t$, through which capitalist accumulation is converted into productive capacity.

\paragraph{B}
Crucially, the comparison between the center (the United States) and the periphery (Chile) is not a comparison between a theoretical norm and a structural deviation. Core and peripheral economies operate within the same fundamental capitalist logic of decentralized accumulation, market anarchy, and internal managerial division. What distinguishes the periphery from the center is not the presence of a different economic logic, but the distinct institutional and structural mediations through which the transformation elasticity $\theta_t$ is determined.

\subsection{Structural Overaccumulation in the Center}

\paragraph{A}
In the central case, capital accumulation proceeds within an advanced, self-reproducing capital-goods sector. The aggregate transformation elasticity $\theta_t$ is determined primarily by domestic class distribution ($\omega$), internal corporate financialization ($\iota$), and the relative asset price of machinery to structures ($p^\kappa$).

\paragraph{B}
Under the post-war Fordist settlement (1929–1973), institutional collective bargaining and strong wage floors generated sustained distributive pressure. In accordance with the Okishio induced-innovation mechanism, cost-minimizing corporate firms responded to wage pressure by accelerating physical mechanization. Because corporate cash flows remained focused on physical plant accumulation rather than financialized asset diversion ($\iota \to 0$), distributive pressure enhanced the capacity-generating efficiency of capital, sustaining high investment dynamism and maintaining capacity utilization within a stable corridor around normal capacity ($\hat{\mu} \in [0.80, 1.05]$).

\paragraph{C}
With the collapse of Fordism and the onset of neoliberal restructuring (1974–2024), this transformation mechanism underwent a profound institutional decoupling. Real wage growth was severed from productivity gains, while corporate liquidity was increasingly diverted into intellectual property products, share buybacks, and financialized asset accumulation ($\iota \uparrow$). Simultaneously, relative machinery price inflation penalized the intensive mechanization margin ($\hat{\theta}_{\kappa p\kappa} < 0$).

\paragraph{D}
This structural shift drove the aggregate transformation elasticity permanently below unity ($\theta_t^{\text{SpecB}} < 1.00$). Bounded between $0.43$ and $0.46$ across the entire post-war era, $\theta_t < 1.00$ defines the \textbf{Structural Overaccumulation Regime}: a historical condition in which every 1\% expansion in physical corporate capital stock yields less than a 1\% expansion in potential capacity output. Under weak labor closure and financialized asset diversion, accumulation decoupling produces chronic industrial underutilization and secular stagnation.

\subsection{The External Realization Trap in the Periphery}

\paragraph{A}
In the peripheral case, the transformation of accumulation into productive capacity encounters a qualitatively distinct structural boundary: the absence of a domestic capital-goods sector. Industrialization under import substitution (ISI) alters the form of external dependence without eliminating it. Instead of importing finished consumer goods, the peripheral economy must import capital equipment ($K^{\text{ME}}$) to expand domestic productive capacity.

\paragraph{B}
Because capital-goods imports must be financed through primary commodity export earnings ($X_t \cdot ToT_t$), the foreign exchange envelope constitutes a binding macroeconomic constraint. Unlike central firms that adjust capital intensity against domestic wage bargaining, peripheral firms confront an external realization bottleneck:

\paragraph{C}
\textbf{Regime A (Slack External Constraint, $D^{ext}_t = 0$).} When the Prebischian Capacity to Import is unconstrained ($C_t > \hat{\gamma}$), domestic wage pressure induces cost-minimizing firms to mechanize ($\hat{\beta}_2 = 4.85^{**}$), replicating the Okishio distributive dividend observed in the core.

\paragraph{D}
\textbf{Regime B (Binding Realization Trap, $D^{ext}_t = 1$).} When terms-of-trade deterioration or export volume contractions depress foreign exchange capacity below the critical threshold ($C_t \le \hat{\gamma}$), the \textbf{Peripheral Realization Penalty} activates violently ($\hat{\delta}_2 = -15.99^{***}$). The net elasticity collapses to $\hat{\beta}_2 + \hat{\delta}_2 = -11.14$.

\paragraph{E}
Under extreme external balance stress, nominal wage increases cannot translate into physical mechanization because foreign exchange scarcity blocks the importation of capital equipment. Surplus value is drained through conspicuous consumption and capital flight, causing capacity utilization to collapse. The balance of payments thus acts as the institutional channel that converts external structural constraints into internal workplace stagnation.

\subsection{Synthesis: Differentiated Determinations of a Common Object}

\paragraph{A}
The comparative findings for the United States and Chile demonstrate that core overaccumulation and peripheral realization traps are mutually entangled expressions of the same world-system reproduction problem. The central economy manages transformation through internal financialization and capital composition shifts, whereas the peripheral economy manages transformation through foreign exchange allocation and external balance constraints.

\paragraph{B}
The difference between center and periphery lies not in whether capacity utilization exists as an identifiable structural object, but in how the transformation of accumulation into productive capacity is structurally mediated across distinct nodes of the world capitalist economy.
\FloatBarrier
\section{Conclusion}

\paragraph{A}
This paper has reconstructed capacity utilization as a structurally mediated object under capitalist disequilibrium. Departing from treatments that view utilization either as an exogenous engineering constant, an unmediated behavioral ratio, or an empirical regression residual, we have demonstrated that capacity utilization ($\mu_t = Y_t / Y_t^p$) becomes identifiable only downstream of the recovery of the aggregate transformation elasticity $\theta_t$.

\paragraph{B}
By integrating Okishio’s theory of decentralized market anarchy, Vidal’s framework of the internally divided firm, and Marglin-style distributive grounding, we established that the transformation of capital accumulation into potential productive capacity ($g_{Y^p} = \theta_t g_K$) is endogenously conditioned by class distribution, corporate financialization, capital composition, and external balance constraints.

\paragraph{C}
Empirically, our cointegrating polynomial regression (CMPR IM-OLS) estimations delivered two major historical conclusions:

\begin{enumerate}
  \item \textbf{Center (United States, 1929–2024):} The U.S. corporate economy operates permanently within a structural regime of capital overaccumulation ($\theta_t < 1.00$). The historical transition from post-war Fordism to neoliberal post-Fordism represents an endogenous regime shift—driven by wage suppression, relative machinery price inflation, and corporate financialization—that generated secular downward drift in capacity utilization.
\end{enumerate}

\begin{enumerate}
  \item \textbf{Periphery (Chile, 1920–2024):} In the peripheral economy, the transformation relation is non-linearly gated by foreign exchange availability. Using Chen’s (2015) Profile Least Squares algorithm over the Prebischian Capacity to Import ($C_t = X_t \cdot ToT_t$), we demonstrated that when external constraints bind ($D^{ext}_t = 1$), the Okishio distributive dividend breaks down, triggering a severe peripheral realization trap ($\hat{\delta}_2 = -15.99^{***}$).
\end{enumerate}

\paragraph{D}
These findings demonstrate that the difference between center and periphery is not a difference between a realized norm and a structural deviation, but a difference in how the same fundamental transformation problem is mediated under distinct world-system conditions.

\paragraph{E}
The scope of this chapter is deliberately bounded. It does not present a complete crisis theory, a fully specified model of profitability dynamics, or an exhaustive dependency framework. The chapter reconstructs capacity utilization as a structurally mediated object and identifies the transformation elasticity as the hinge of its empirical recovery across center and periphery, while leaving the broader crisis and profitability architecture for subsequent work.

% ---------- Appendix ----------%
% \appendix precedes the bibliography so appendix sections carry their own lettered
% numbering instead of being absorbed into the preceding section's outline branch.
\pagebreak
\appendix
\FloatBarrier
\appendix
\section*{Appendix A: Proof of Proposition 3}

\paragraph{A.1}
This appendix provides the detailed derivation of the
comparative statics stated in Proposition~3
(\S\ref{eq:32:cs_omega}--\ref{eq:32:cs_price}). The proof
proceeds by differentiating the optimal mechanization gap
\eqref{eq:32:kappa_star} with respect to each state
variable.

\subsection*{A.1 The Optimal Mechanization Gap}

\paragraph{A.2}
Recall from \S3.2 that the firm maximizes the net
viability surplus:
\begin{equation}
	\max_{\kappa_t}\; c(\kappa_t)
	= \Phi(\kappa_t) - \kappa_t \cdot MC(\omega_{t-1}, p^\kappa_t)
	\label{eq:A:objective}
\end{equation}
where $\Phi(\kappa_t) = \lambda_0 + \lambda_1 \kappa_t
+ \lambda_2 \kappa_t^2$ is the mechanization frontier
\eqref{eq:32:frontier_kappa}, and
$MC(\omega_{t-1}, p^\kappa_t) = \gamma_0 - \gamma_1 \omega_{t-1}
+ \gamma_2 p^\kappa_t$ is the marginal viability cost
\eqref{eq:32:MC}. Substituting and differentiating:
\begin{align}
	c(\kappa_t) &= \lambda_0
	+ \left[\lambda_1 - \gamma_0 + \gamma_1 \omega_{t-1}
	- \gamma_2 p^\kappa_t\right] \kappa_t
	+ \lambda_2 \kappa_t^2
	\label{eq:A:unconstrained} \\
	\frac{dc}{d\kappa_t} &= \lambda_1 - \gamma_0
	+ \gamma_1 \omega_{t-1} - \gamma_2 p^\kappa_t
	+ 2\lambda_2 \kappa_t^* = 0
	\label{eq:A:foc}
\end{align}
Solving \eqref{eq:A:foc} under Assumption~1
($\lambda_2 < 0$) yields the optimal mechanization gap:
\begin{equation}
	\kappa_t^*(\omega_{t-1}, p^\kappa_t)
	= \frac{(\gamma_0 - \lambda_1) - \gamma_1 \omega_{t-1}
		+ \gamma_2 p^\kappa_t}{2\lambda_2}
	\label{eq:A:kappa_star}
\end{equation}
which reproduces \eqref{eq:32:kappa_star}.

\subsection*{A.2 The Distributive Channel}

\paragraph{A.3}
\textbf{Proof of} \eqref{eq:32:cs_omega}. Differentiating
\eqref{eq:A:kappa_star} with respect to $\omega_{t-1}$:
\begin{equation}
	\frac{\partial \kappa_t^*}{\partial \omega_{t-1}}
	= \frac{-\gamma_1}{2\lambda_2}
	\label{eq:A:cs_omega}
\end{equation}
Under Assumption~1, $\lambda_2 < 0$. By Assumption~3
($\gamma_1 > 0$), the numerator $-\gamma_1 < 0$ and the
denominator $2\lambda_2 < 0$. Therefore:
\begin{equation}
	\frac{\partial \kappa_t^*}{\partial \omega_{t-1}}
	= \frac{-\gamma_1}{2\lambda_2} > 0
	\label{eq:A:cs_omega_sign}
\end{equation}
A rising wage share forces cost-minimizing firms to widen
the mechanization gap, substituting machinery for labor.
$\hfill\blacksquare$

\subsection*{A.3 The Relative Price Channel}

\paragraph{A.4}
\textbf{Proof of} \eqref{eq:32:cs_price}. Differentiating
\eqref{eq:A:kappa_star} with respect to $p^\kappa_t$:
\begin{equation}
	\frac{\partial \kappa_t^*}{\partial p^\kappa_t}
	= \frac{\gamma_2}{2\lambda_2}
	\label{eq:A:cs_price}
\end{equation}
Under Assumption~1, $\lambda_2 < 0$. By Assumption~3
($\gamma_2 > 0$), the numerator $\gamma_2 > 0$ and the
denominator $2\lambda_2 < 0$. Therefore:
\begin{equation}
	\frac{\partial \kappa_t^*}{\partial p^\kappa_t}
	= \frac{\gamma_2}{2\lambda_2} < 0
	\label{eq:A:cs_price_sign}
\end{equation}
A rising relative price of machinery suppresses the optimal
mechanization gap, as the structural cost of the composition
shift penalizes mechanization. $\hfill\blacksquare$

\subsection*{A.4 The Cost Hurdle Formulation (Alternative
	Derivation)}

\paragraph{A.5}
An equivalent derivation expresses the FOC as a cost
hurdle. Define the effective cost hurdle of mechanization:
\begin{equation}
	H(\omega_{t-1}, p^\kappa_t)
	\equiv \frac{MC(\omega_{t-1}, p^\kappa_t)}{1}
	= \gamma_0 - \gamma_1 \omega_{t-1} + \gamma_2 p^\kappa_t
	\label{eq:A:hurdle}
\end{equation}
The FOC \eqref{eq:A:foc} can be rewritten as:
\begin{equation}
	\Phi'(\kappa_t^*) = H(\omega_{t-1}, p^\kappa_t)
	\label{eq:A:foc_hurdle}
\end{equation}
The firm mechanizes up to the point where the marginal
productivity of the mechanization gap equals the effective
cost hurdle. The comparative statics follow from the
implicit function theorem applied to
\eqref{eq:A:foc_hurdle}:
\begin{align}
	\frac{\partial \kappa_t^*}{\partial \omega_{t-1}}
	&= -\frac{\partial H / \partial \omega_{t-1}}
	{\Phi''(\kappa_t^*)}
	= -\frac{-\gamma_1}{2\lambda_2} > 0
	\label{eq:A:ift_omega} \\
	\frac{\partial \kappa_t^*}{\partial p^\kappa_t}
	&= -\frac{\partial H / \partial p^\kappa_t}
	{\Phi''(\kappa_t^*)}
	= -\frac{\gamma_2}{2\lambda_2} < 0
	\label{eq:A:ift_price}
\end{align}
Both derivations yield identical results, confirming
Proposition~3.
\FloatBarrier
\section{Data Construction and Measurement Provenance}

\paragraph{A}
The main text of §4.1 keeps only what identification requires: the constructed capital stocks, the distinction between depreciation and retirement, the analytical roles of gross and net capital, and the rule that cross-country comparability rests on accounting consistency rather than on harmonized source vintages. Everything needed to reproduce those constructions, but not to follow the argument, is collected here.

\subsection{The generalized perpetual inventory method}

\paragraph{A}
\PLACEHOLDER{Full statement of the GPIM construction: survival function and its parameterization, service-life assumptions by asset class, the initialization of the capital stock at the start of each series, and the treatment of the discard rate MATHPLACEHOLDER0XYZ on gross capital in operation as distinct from the net depreciation rate MATHPLACEHOLDER1XYZ. Source material sits in the archived maintenance audits at \texttt{reports/maintenance/} --- specifically the gross-survival remediation, input-initialization, price-deflator provenance, and source-price-seed audits.}

\subsection{Price bases and deflator provenance}

\paragraph{A}
\PLACEHOLDER{Deflator chain for each series and the base-year conventions. The Chilean capital stock and GDP series share a 2003 CLP base and require no rebasing, which is the reason the Chilean accounting closes without a splice at that seam; this needs stating explicitly with its source. Pre-1940 Chilean relative capital-asset prices are proxied by a US wholesale price index, and the consequences of that proxy for the price channel need to be stated here rather than left in a table footnote.}

\subsection{Source reconciliation and splice diagnostics}

\paragraph{A}
\PLACEHOLDER{Series-by-series reconciliation against source vintages, splice points and the diagnostics run at each, and the repository provenance of every input. Source material is in the archived \texttt{reports/US\_Wage\_Share\_Data\_Provenance\_Note/99\_archive/} and the corresponding Chilean lane.}

\subsection{The wage-share measures}

\paragraph{A}
Four wage-share measures enter the estimation and the distinctions between them do analytical work, so they are set out here in full rather than by reference.

\paragraph{B}
The baseline measure takes compensation over value added for the non-financial corporate sector. The gross and net variants differ in whether consumption of fixed capital sits in the denominator; because that quantity rises mechanically with capital intensity, the two variants respond differently to exactly the mechanization the interaction term is meant to identify, which is why §4.2 treats the gross-basis failure as evidence rather than as noise.

\paragraph{C}
The two Shaikh-Tonak variants further separate productive from supervisory labour compensation. Tier 1 isolates variable capital on a first pass; Tier 2 applies the supervisory correction in full. The theoretical motivation is that cost-minimizing choice of technique responds to shop-floor labour cost, not to an aggregate that folds managerial overhead into the same ratio.

\paragraph{D}
\PLACEHOLDER{Construction detail for each of the four measures --- the exact NIPA line items, the supervisory split and its source, and the crosswalk between the four --- belongs in this subsection.}

\subsection{The level anchor}

\paragraph{A}
Capacity utilization is a ratio whose numerator is observed and whose denominator is estimated. Estimation recovers the transformation relation and therefore the shape of the productive-capacity path, but not its level: any scalar multiple of $\hat{Y}^p_t$ fits equally well. The level is fixed by anchoring utilization to unity in a benchmark year chosen to approximate frontier utilization.

\paragraph{B}
This normalization anchors the series. It does not define the object. The distinction matters because it is precisely what separates this construction from approaches that recover utilization as a residual: here the transformation elasticity is identified first and the level is imposed afterwards, so the anchor year carries no weight in identifying $\theta$.

\paragraph{C}
\PLACEHOLDER{Statement of the pinch-year selection rule, the chosen benchmark year for each case with its justification, and sensitivity of the reconstructed level path to alternative anchor years. Protocol is at \texttt{chapter2\_vault/04\_data\_measurement/D03\_capacity\_utilization\_level\_anchor\_pinch\_year\_protocol.md}.}
\FloatBarrier
\section{The Specification Space and Model Selection}

\paragraph{A}
The main text reports one preferred specification per case. This appendix reports the space those specifications were selected from, which alternatives were estimated, and why each is not carried forward. Presenting a single preferred model without its search space would leave the reader unable to judge how much the result depends on choices made along the way.

\subsection{The combinatorial grid}

\paragraph{A}
Specification A varies three binary dimensions — wage-share denominator (gross versus net value added), interaction form (standard versus effective, the latter scaled by the financialization retention index $(1-\iota_{t-1})$), and the presence of Bai-Perron step level-shifts — across three wage-share measures, giving the estimation matrix reported in §4.2. Specification B adds a fourth dimension by disaggregating capital into structures $k^{\text{NRC}}$ and the mechanization gap $\kappa_t$, and by admitting a relative-price interaction alongside the distributive one.

\paragraph{B}
Every cell of these grids was estimated. None was dropped for producing an inconvenient result.

\subsection{Why alternatives are not carried forward}

\paragraph{A}
Specifications fall out of the preferred set for reasons of different weight, and the distinction matters for interpretation.

\paragraph{B}
Some are excluded on \textbf{admissibility}: a regressor fails the integration-order gate as actually constructed, including its lag structure, or a generated interaction inherits $I(2)$ behaviour from its components. Theoretical motivation does not rescue these. The theory describes a decision mechanism; the stochastic evidence determines whether a given observable proxy for that mechanism can enter a long-run regression at all.

\paragraph{C}
Some are excluded on \textbf{measurement}: the gross-value-added wage shares are admissible but uninformative, because consumption of fixed capital rises mechanically with capital intensity and so contaminates the distributive signal precisely along the dimension the interaction is meant to identify. Their failure is itself a result, not a nuisance, and §4.2 reads it as one.

\paragraph{D}
Some are excluded on \textbf{redundancy}: the step-dummy variants sustain cointegration but add nothing once the distributive-financial interaction is specified, which is the basis for carrying forward the no-break specification.

\paragraph{E}
None of the estimated specifications was excluded because it failed to cointegrate.

\paragraph{F}
\PLACEHOLDER{The full consolidated grid --- the MATHPLACEHOLDER0XYZ Specification A and MATHPLACEHOLDER1XYZ Specification B permutations under initial-state centering, with each cell's disposition and the economic reason for it --- is already written at \texttt{chapter2\_vault/06\_paper\_facing/results/99\_legacy/S05\_Appendix.md} (status: locked). It needs transposing into this appendix with its internal governance vocabulary translated into plain econometric language: the vault's four-way ledger of dispositions is an internal audit device and must not surface in paper-facing text.}

\paragraph{G}
\PLACEHOLDER{The integration-order verdict table --- each candidate regressor mapped to its theoretical basis, its MATHPLACEHOLDER0XYZ evidence on both levels and first differences from at least two complementary tests, and its final disposition --- belongs here. Source material is in the archived \texttt{reports/US\_Integration\_Order\_Appendix/99\_archive/}.}

\subsection{Counterfactual restrictions}

\paragraph{A}
Two restrictions test whether the estimated transformation elasticity is doing work the data support rather than work the specification imposes.

\paragraph{B}
The first restricts $\theta = 1$, the Harrodian knife-edge, and asks what the capacity series and the implied utilization path look like under that restriction. This is the case in which accumulation converts one-for-one into productive capacity, and it is the implicit assumption of any filtering approach that fits output to capital without separately identifying the transformation relation.

\paragraph{C}
The second removes the step dummies entirely and re-estimates, to establish how much of the reported structure survives without any exogenous break structure.

\paragraph{D}
\PLACEHOLDER{Both counterfactuals are specified but not yet reported. Each requires a re-estimation pass and a comparison table against the preferred specification.}
\FloatBarrier
\section{Cointegration Diagnostics and the Monte Carlo Critical-Value Protocol}

\paragraph{A}
The main text reports one cointegration statistic per specification: the Vogelsang and Wagner (2014) Type 2 test. This appendix reports the full diagnostic set behind that verdict, states the critical values, and explains why the Type 2 test rather than the residual tests carries the reading.

\subsection{Why the Type 2 test governs}

\paragraph{A}
The three tests in use do not stand on equal footing, and treating them as a co-equal triad obscures what each is doing.

\paragraph{B}
The residual Engle-Granger ADF test asks whether the estimated residual series contains a unit root. Its null is \textit{no} cointegration, so rejection supports the long-run relation. The Shin LM test inverts this: its null is cointegration, and rejection indicates spurious regression. The Vogelsang-Wagner Type 2 test also takes cointegration as its null, but it differs from both residual tests in a way that matters for this chapter. Rather than testing a property of a residual series constructed after the fact, it tests a restriction on the IM-OLS estimator itself — that the auxiliary coefficient vector $\lambda_2$ is zero, which holds under genuine cointegration and fails under spurious regression. Because the specifications here carry generated interaction regressors that have been FWL-orthogonalized against the base capital terms and the deterministic space, a residual-based test inherits whatever the orthogonalization did to the residual's serial dependence. The Type 2 statistic does not.

\paragraph{C}
That is the substantive reason for the hierarchy, and it is the one the estimation protocol adopts: the Type 2 test is primary and binding, the residual ADF and Shin LM tests are secondary robustness checks, and Johansen and Phillips-Ouliaris statistics are background diagnostics that do not enter the paper.

\subsection{Reading the stars: opposed nulls}

\paragraph{A}
Because the nulls run in opposite directions, an identical typographic convention carries opposite meanings across rows of the tables below. This is worth stating plainly rather than leaving to a footnote.

\paragraph{B}
For the residual ADF test, stars indicate rejection of no cointegration, and are therefore the favourable result. For the Shin LM and Vogelsang-Wagner Type 2 tests, stars indicate rejection of cointegration, and are therefore the unfavourable result. A well-behaved specification shows stars on its ADF row and none on the other two.

\subsection{Critical values}

\begin{table}[H]
\centering
\caption{Critical Values for the Cointegration Diagnostics}
\label{tab:critical_values_for_the_cointegration_diagnostic}
\begin{tabular}{p{5.5cm}p{2.2cm}p{2.2cm}p{2.2cm}p{2.2cm}}
\toprule
Test & Source & 10\% & 5\% & 1\% \\
\midrule
Residual ADF ($t_{\text{ADF}}$) & MacKinnon (2010), Case 2, constant only, $m=4$ & $-3.81$ & $-4.10$ & $-4.64$ \\
Shin LM ($C_{\text{Shin}}$), $k=2$ & Shin (1994), Table 1, demeaned $C_\mu$ & $0.163$ & $0.221$ & $0.380$ \\
Shin LM ($C_{\text{Shin}}$), $k=3$ & Shin (1994), Table 1, demeaned $C_\mu$ & $0.121$ & $0.159$ & $0.271$ \\
VW Type 2 ($S_{\text{Type 2}}$) & Vogelsang and Wagner (2014), fixed-$b$, Bartlett kernel, intercept-only & jointly indexed by $(b,k)$ & no single constant applies & see D.3 note \\
\bottomrule
\end{tabular}
\end{table}

\paragraph{A}
Here $m$ counts the total number of $I(1)$ variables in the system, one dependent variable plus the stochastic regressors, with step dummies excluded from the count; $k$ counts stochastic $I(1)$ regressors alone. Exogenous $I(0)$ controls sit in the deterministic matrix and so do not enter $k$.

\paragraph{B}
The last row is the substantive point. Vogelsang-Wagner fixed-$b$ critical values are not constants. They vary by roughly an order of magnitude across the tabulated bandwidth-ratio range and depend on the regressor count, so every model must be evaluated against the cell matching its own $(b, k)$ pair. Each table below reports the bandwidth ratio $b$ for exactly this reason.

\paragraph{C}
\PLACEHOLDER{The Type 2 rows in the tables below carry no significance stars, because the statistics have not yet been evaluated against the MATHPLACEHOLDER0XYZ-indexed fixed-MATHPLACEHOLDER1XYZ table digitized in \texttt{codes/COMMON\_VW\_fixedb\_critical\_values.R}. Earlier drafts compared them against a single hardcoded triple (MATHPLACEHOLDER2XYZ / MATHPLACEHOLDER3XYZ / MATHPLACEHOLDER4XYZ) applied uniformly across models regardless of bandwidth or regressor count; a 2026-08-03 audit established that this triple matches no cell of the published table, and those comparisons have been withdrawn rather than reproduced here. Re-running the estimation against the digitized module is the outstanding work. Until it is done, this appendix reports the statistics and the critical-value regime, but not a starred verdict.}

\subsection{Constructing the residual statistics}

\paragraph{A}
Two implementation details determine whether the residual diagnostics measure what they are supposed to measure.

\paragraph{B}
The Shin LM statistic is evaluated on the IM-OLS partial sum residuals $U_t = \sum e_t$ directly, with no further accumulation. Because $U_t$ is already an integrated partial sum by construction, applying a second cumulative sum produces an $I(2)$ series and inflates the statistic. The error is easy to make and hard to see, since the inflated statistic remains finite and ordered plausibly across models.

\paragraph{C}
The regressors entering the residual tests are the deterministically detrended series — the residuals of $\mathbf{x}$ regressed on the step dummies and their partial sums — from which the reverse cumulative sum regressors are then built. Skipping the detrending leaves deterministic trend contamination in the test.

\subsection{Small-sample and lag-augmentation caveats}

\paragraph{A}
Two further constraints bear on how the residual diagnostics below should be read.

\paragraph{B}
Residual unit-root tests must be lag-augmented with an information-criterion-selected lag order. A plain Dickey-Fuller regression with zero lagged differences is not admissible where the residuals are serially correlated, which in this setting means everywhere. Where a sub-sample is short — the Fordist window gives $T=44$ — small-sample bias pushes the true critical values more negative than the asymptotic boundaries tabulated above, so marginal significance near $-4.10$ in a sub-sample should not be read as a clean rejection.

\paragraph{C}
Second, step dummies test only for level shifts. They cannot absorb a shift in the cointegrating vector itself, so a specification that answers detected breaks with step dummies alone has not addressed slope instability and is not described as break-robust anywhere in this chapter. Appendix E takes up parameter stability directly.

\subsection{Specification A: aggregate capital}

\begin{table}[H]
\centering
\setlength{\tabcolsep}{2.5pt}
\caption{Specification Set 1: Standard Wage Interaction (GVA Basis) --- Cointegration Diagnostics and Collinearity}
\label{tab:specA_set1_gva_triad}
\scriptsize
\begin{tabular}{p{3.8cm}p{1.1cm}p{1.1cm}p{1.1cm}p{1.1cm}p{1.1cm}p{1.1cm}p{1.1cm}p{1.1cm}}
\toprule
\multicolumn{9}{l}{\textbf{Dependent Variable: $y_t^{\text{NFC}}$ (Log NFC Real NVA)}} \\
\midrule
Model & M1.1a & M1.1b & M1.2a & M1.2b & M1.3a & M1.3b & M1.4a & M1.4b \\
\midrule
\multicolumn{9}{l}{\textbf{Collinearity Diagnostics (VIF)}} \\
VIF ($k_{\text{cap}}$) & 1.00 & 1.00 & 1.00 & 1.00 & 1.00 & 1.00 & 1.00 & 1.00 \\
VIF (Interaction) & 1.00 & 1.00 & 1.00 & 1.00 & 1.00 & 1.00 & 1.00 & 1.00 \\
\midrule
\multicolumn{9}{l}{\textbf{Cointegration Test Triad}} \\
Resid. ADF ($t_{\text{ADF}}$) [$H_0$: No Coint.] & -2.41 & -2.81 & -2.33 & -2.73 & -4.15$^{**}$ & -3.56$^{*}$ & -3.40 & -3.55$^{*}$ \\
Shin LM ($C_{\text{Shin}}$) [$H_0$: Coint.] & 0.06 & 0.04 & 0.06 & 0.04 & 0.05 & 0.03 & 0.03 & 0.03 \\
VW Type 2 ($S_{\text{Type 2}}$) [$H_0$: Coint.] & 3.30 & 11.34 & 2.70 & 11.41 & 5.82 & 9.68 & 1.46 & 9.08 \\
\midrule
Bandwidth Ratio ($b$) & 0.2746 & 0.1872 & 0.2864 & 0.1934 & 0.2852 & 0.1997 & 0.2665 & 0.2731 \\
Sample Size ($T$) & 94 & 94 & 94 & 94 & 94 & 94 & 94 & 94 \\
\bottomrule
\multicolumn{9}{p{14.5cm}}{\footnotesize \textbf{Note:} Full cointegration diagnostics for the estimates reported in the main text. \textbf{Star reversal warning:} for the Residual ADF test, significance stars indicate rejection of No Cointegration (supporting cointegration). For the Shin LM test, significance stars indicate rejection of Cointegration (indicating spurious regression). The Vogelsang--Wagner Type 2 test shares the Shin null ($H_0$: cointegration) and its stars read the same way. Critical values: Residual ADF (MacKinnon 2010, Case 2, constant only, $m=4$) $-3.81$ / $-4.10$ / $-4.64$ at 10\% / 5\% / 1\%; Shin LM (Shin 1994, Table 1, demeaned $C_\mu$) $0.163$ / $0.221$ / $0.380$ for $k=2$ and $0.121$ / $0.159$ / $0.271$ for $k=3$, where $k$ counts stochastic $I(1)$ regressors; Vogelsang--Wagner Type 2 critical values are fixed-$b$ and jointly indexed by $(b, k)$, so no single constant applies across models \textbf{[PENDING: see the note to the corresponding main-text table]}. Bandwidth ratio $b$ is the model's HAC bandwidth divided by $T$.} \\
\end{tabular}
\end{table}


\begin{table}[H]
\centering
\setlength{\tabcolsep}{2.5pt}
\caption{Specification Set 1: Standard Wage Interaction (NVA Basis) --- Cointegration Diagnostics and Collinearity}
\label{tab:specA_set1_nva_triad}
\scriptsize
\begin{tabular}{p{3.8cm}p{1.1cm}p{1.1cm}p{1.1cm}p{1.1cm}p{1.1cm}p{1.1cm}p{1.1cm}p{1.1cm}}
\toprule
\multicolumn{9}{l}{\textbf{Dependent Variable: $y_t^{\text{NFC}}$ (Log NFC Real NVA)}} \\
\midrule
Model & M1.5a & M1.5b & M1.6a & M1.6b & M1.7a & M1.7b & M1.8a & M1.8b \\
\midrule
\multicolumn{9}{l}{\textbf{Collinearity Diagnostics (VIF)}} \\
VIF ($k_{\text{cap}}$) & 1.00 & 1.00 & 1.00 & 1.00 & 1.00 & 1.00 & 1.00 & 1.00 \\
VIF (Interaction) & 1.00 & 1.00 & 1.00 & 1.00 & 1.00 & 1.00 & 1.00 & 1.00 \\
\midrule
\multicolumn{9}{l}{\textbf{Cointegration Test Triad}} \\
Resid. ADF ($t_{\text{ADF}}$) [$H_0$: No Coint.] & -2.41 & -2.81 & -2.41 & -2.81 & -2.41 & -2.81 & -2.94 & -2.95 \\
Shin LM ($C_{\text{Shin}}$) [$H_0$: Coint.] & 0.06 & 0.04 & 0.06 & 0.04 & 0.06 & 0.04 & 0.04 & 0.03 \\
VW Type 2 ($S_{\text{Type 2}}$) [$H_0$: Coint.] & 3.30 & 11.34 & 3.30 & 11.34 & 3.30 & 11.34 & 0.05 & 11.94 \\
\midrule
Bandwidth Ratio ($b$) & 0.2746 & 0.1872 & 0.2746 & 0.1872 & 0.2746 & 0.1872 & 0.1940 & 0.1935 \\
Sample Size ($T$) & 94 & 94 & 94 & 94 & 94 & 94 & 94 & 94 \\
\bottomrule
\multicolumn{9}{p{14.5cm}}{\footnotesize \textbf{Note:} Full cointegration diagnostics for the estimates reported in the main text. \textbf{Star reversal warning:} for the Residual ADF test, significance stars indicate rejection of No Cointegration (supporting cointegration). For the Shin LM test, significance stars indicate rejection of Cointegration (indicating spurious regression). The Vogelsang--Wagner Type 2 test shares the Shin null ($H_0$: cointegration) and its stars read the same way. Critical values: Residual ADF (MacKinnon 2010, Case 2, constant only, $m=4$) $-3.81$ / $-4.10$ / $-4.64$ at 10\% / 5\% / 1\%; Shin LM (Shin 1994, Table 1, demeaned $C_\mu$) $0.163$ / $0.221$ / $0.380$ for $k=2$ and $0.121$ / $0.159$ / $0.271$ for $k=3$, where $k$ counts stochastic $I(1)$ regressors; Vogelsang--Wagner Type 2 critical values are fixed-$b$ and jointly indexed by $(b, k)$, so no single constant applies across models \textbf{[PENDING: see the note to the corresponding main-text table]}. Bandwidth ratio $b$ is the model's HAC bandwidth divided by $T$.} \\
\end{tabular}
\end{table}


\begin{table}[H]
\centering
\setlength{\tabcolsep}{2.5pt}
\caption{Specification Set 2: Effective Wage Interaction (GVA Basis) --- Cointegration Diagnostics and Collinearity}
\label{tab:specA_set2_gva_triad}
\scriptsize
\begin{tabular}{p{3.8cm}p{1.1cm}p{1.1cm}p{1.1cm}p{1.1cm}p{1.1cm}p{1.1cm}p{1.1cm}p{1.1cm}}
\toprule
\multicolumn{9}{l}{\textbf{Dependent Variable: $y_t^{\text{NFC}}$ (Log NFC Real NVA)}} \\
\midrule
Model & M2.1a & M2.1b & M2.2a & M2.2b & M2.3a & M2.3b & M2.4a & M2.4b \\
\midrule
\multicolumn{9}{l}{\textbf{Collinearity Diagnostics (VIF)}} \\
VIF ($k_{\text{cap}}$) & 1.00 & 1.00 & 1.00 & 1.00 & 1.00 & 1.00 & 1.00 & 1.00 \\
VIF (Interaction) & 1.00 & 1.00 & 1.00 & 1.00 & 1.00 & 1.00 & 1.00 & 1.00 \\
\midrule
\multicolumn{9}{l}{\textbf{Cointegration Test Triad}} \\
Resid. ADF ($t_{\text{ADF}}$) [$H_0$: No Coint.] & -2.32 & -3.06 & -2.26 & -2.96 & -4.02$^{**}$ & -2.93 & -3.27 & -2.90 \\
Shin LM ($C_{\text{Shin}}$) [$H_0$: Coint.] & 0.06 & 0.04 & 0.06 & 0.04 & 0.05 & 0.03 & 0.03 & 0.03 \\
VW Type 2 ($S_{\text{Type 2}}$) [$H_0$: Coint.] & 10.76 & 9.07 & 6.17 & 9.61 & 4.39 & 7.88 & 0.73 & 4.55 \\
\midrule
Bandwidth Ratio ($b$) & 0.2976 & 0.1687 & 0.3066 & 0.1732 & 0.2908 & 0.1765 & 0.2713 & 0.2047 \\
Sample Size ($T$) & 94 & 94 & 94 & 94 & 94 & 94 & 94 & 94 \\
\bottomrule
\multicolumn{9}{p{14.5cm}}{\footnotesize \textbf{Note:} Full cointegration diagnostics for the estimates reported in the main text. \textbf{Star reversal warning:} for the Residual ADF test, significance stars indicate rejection of No Cointegration (supporting cointegration). For the Shin LM test, significance stars indicate rejection of Cointegration (indicating spurious regression). The Vogelsang--Wagner Type 2 test shares the Shin null ($H_0$: cointegration) and its stars read the same way. Critical values: Residual ADF (MacKinnon 2010, Case 2, constant only, $m=4$) $-3.81$ / $-4.10$ / $-4.64$ at 10\% / 5\% / 1\%; Shin LM (Shin 1994, Table 1, demeaned $C_\mu$) $0.163$ / $0.221$ / $0.380$ for $k=2$ and $0.121$ / $0.159$ / $0.271$ for $k=3$, where $k$ counts stochastic $I(1)$ regressors; Vogelsang--Wagner Type 2 critical values are fixed-$b$ and jointly indexed by $(b, k)$, so no single constant applies across models \textbf{[PENDING: see the note to the corresponding main-text table]}. Bandwidth ratio $b$ is the model's HAC bandwidth divided by $T$.} \\
\end{tabular}
\end{table}


\begin{table}[H]
\centering
\setlength{\tabcolsep}{2.5pt}
\caption{Specification Set 2: Effective Wage Interaction (NVA Basis) --- Cointegration Diagnostics and Collinearity}
\label{tab:specA_set2_nva_triad}
\scriptsize
\begin{tabular}{p{3.8cm}p{1.1cm}p{1.1cm}p{1.1cm}p{1.1cm}p{1.1cm}p{1.1cm}p{1.1cm}p{1.1cm}}
\toprule
\multicolumn{9}{l}{\textbf{Dependent Variable: $y_t^{\text{NFC}}$ (Log NFC Real NVA)}} \\
\midrule
Model & M2.5a & M2.5b & M2.6a & M2.6b & M2.7a & M2.7b & M2.8a & M2.8b \\
\midrule
\multicolumn{9}{l}{\textbf{Collinearity Diagnostics (VIF)}} \\
VIF ($k_{\text{cap}}$) & 1.00 & 1.00 & 1.00 & 1.00 & 1.00 & 1.00 & 1.00 & 1.00 \\
VIF (Interaction) & 1.00 & 1.00 & 1.00 & 1.00 & 1.00 & 1.00 & 1.00 & 1.00 \\
\midrule
\multicolumn{9}{l}{\textbf{Cointegration Test Triad}} \\
Resid. ADF ($t_{\text{ADF}}$) [$H_0$: No Coint.] & -2.32 & -3.06 & -2.32 & -3.06 & -2.32 & -3.06 & -2.96 & -3.82$^{**}$ \\
Shin LM ($C_{\text{Shin}}$) [$H_0$: Coint.] & 0.06 & 0.04 & 0.06 & 0.04 & 0.06 & 0.04 & 0.04 & 0.03 \\
VW Type 2 ($S_{\text{Type 2}}$) [$H_0$: Coint.] & 10.76 & 9.07 & 10.76 & 9.07 & 10.76 & 9.07 & 0.46 & 13.69 \\
\midrule
Bandwidth Ratio ($b$) & 0.2976 & 0.1687 & 0.2976 & 0.1687 & 0.2976 & 0.1687 & 0.1873 & 0.1651 \\
Sample Size ($T$) & 94 & 94 & 94 & 94 & 94 & 94 & 94 & 94 \\
\bottomrule
\multicolumn{9}{p{14.5cm}}{\footnotesize \textbf{Note:} Full cointegration diagnostics for the estimates reported in the main text. \textbf{Star reversal warning:} for the Residual ADF test, significance stars indicate rejection of No Cointegration (supporting cointegration). For the Shin LM test, significance stars indicate rejection of Cointegration (indicating spurious regression). The Vogelsang--Wagner Type 2 test shares the Shin null ($H_0$: cointegration) and its stars read the same way. Critical values: Residual ADF (MacKinnon 2010, Case 2, constant only, $m=4$) $-3.81$ / $-4.10$ / $-4.64$ at 10\% / 5\% / 1\%; Shin LM (Shin 1994, Table 1, demeaned $C_\mu$) $0.163$ / $0.221$ / $0.380$ for $k=2$ and $0.121$ / $0.159$ / $0.271$ for $k=3$, where $k$ counts stochastic $I(1)$ regressors; Vogelsang--Wagner Type 2 critical values are fixed-$b$ and jointly indexed by $(b, k)$, so no single constant applies across models \textbf{[PENDING: see the note to the corresponding main-text table]}. Bandwidth ratio $b$ is the model's HAC bandwidth divided by $T$.} \\
\end{tabular}
\end{table}


\subsection{Specification B: heterogeneous capital}

\begin{table}[H]
\centering
\setlength{\tabcolsep}{2.5pt}
\caption{Specification Set 1: Standard Wage Interaction (GVA Basis) --- Cointegration Diagnostics and Collinearity}
\label{tab:specB_set1_gva_triad}
\scriptsize
\begin{tabular}{p{3.8cm}p{1.1cm}p{1.1cm}p{1.1cm}p{1.1cm}p{1.1cm}p{1.1cm}}
\toprule
\multicolumn{7}{l}{\textbf{Dependent Variable: $y_t^{\text{NFC}}$ (Log NFC Real GVA)}} \\
\midrule
Model & M2.1a & M2.2b & M2.5a & M2.6b & M2.9a & M2.10b \\
\midrule
\multicolumn{7}{l}{\textbf{Collinearity Diagnostics (VIF)}} \\
VIF ($k^{\text{NRC}}$) & 1.69 & 1.69 & 1.69 & 1.69 & 1.69 & 1.69 \\
VIF ($\kappa$) & 1.69 & 1.69 & 1.69 & 1.69 & 1.69 & 1.69 \\
VIF (Wage Interaction) & 1.32 & 1.06 & 1.15 & 1.01 & 1.01 & 1.01 \\
VIF (Price Interaction) & 1.32 & 1.06 & 1.15 & 1.01 & 1.01 & 1.01 \\
\midrule
\multicolumn{7}{l}{\textbf{Cointegration Test Triad}} \\
Resid. ADF ($t_{\text{ADF}}$) [$H_0$: No Coint.] & -4.16$^{*}$ & -4.46$^{**}$ & -4.26$^{*}$ & -4.69$^{**}$ & -4.45$^{**}$ & -4.80$^{**}$ \\
Shin LM ($C_{\text{Shin}}$) [$H_0$: Coint.] & 0.05 & 0.03 & 0.05 & 0.03 & 0.04 & 0.02 \\
VW Type 2 ($S_{\text{Type 2}}$) [$H_0$: Coint.] & 16.11 & 9.55 & 16.25 & 11.41 & 16.74 & 7.66 \\
\midrule
Bandwidth Ratio ($b$) & 0.2088 & 0.1758 & 0.2114 & 0.1744 & 0.2088 & 0.1759 \\
Sample Size ($T$) & 94 & 94 & 94 & 94 & 94 & 94 \\
\bottomrule
\multicolumn{7}{p{15cm}}{\footnotesize \textbf{Note:} Full cointegration diagnostics for the estimates reported in the main text. \textbf{Star reversal warning:} for the Residual ADF test, significance stars indicate rejection of No Cointegration (supporting cointegration). For the Shin LM test, significance stars indicate rejection of Cointegration (indicating spurious regression). The Vogelsang--Wagner Type 2 test shares the Shin null ($H_0$: cointegration) and its stars read the same way. Critical values: Residual ADF (MacKinnon 2010, Case 2, constant only, $m=4$) $-3.81$ / $-4.10$ / $-4.64$ at 10\% / 5\% / 1\%; Shin LM (Shin 1994, Table 1, demeaned $C_\mu$) $0.163$ / $0.221$ / $0.380$ for $k=2$ and $0.121$ / $0.159$ / $0.271$ for $k=3$, where $k$ counts stochastic $I(1)$ regressors; Vogelsang--Wagner Type 2 critical values are fixed-$b$ and jointly indexed by $(b, k)$, so no single constant applies across models \textbf{[PENDING: see the note to the corresponding main-text table]}. Bandwidth ratio $b$ is the model's HAC bandwidth divided by $T$.} \\
\end{tabular}
\end{table}


\begin{table}[H]
\centering
\setlength{\tabcolsep}{2.5pt}
\caption{Specification Set 1: Standard Wage Interaction (NVA Basis) --- Cointegration Diagnostics and Collinearity}
\label{tab:specB_set1_nva_triad}
\scriptsize
\begin{tabular}{p{3.8cm}p{1.1cm}p{1.1cm}p{1.1cm}p{1.1cm}p{1.1cm}p{1.1cm}}
\toprule
\multicolumn{7}{l}{\textbf{Dependent Variable: $y_t^{\text{NFC}}$ (Log NFC Real GVA)}} \\
\midrule
Model & M2.3a & M2.4b & M2.7a & M2.8b & M2.11a & M2.12b \\
\midrule
\multicolumn{7}{l}{\textbf{Collinearity Diagnostics (VIF)}} \\
VIF ($k^{\text{NRC}}$) & 1.69 & 1.69 & 1.69 & 1.69 & 1.69 & 1.69 \\
VIF ($\kappa$) & 1.69 & 1.69 & 1.69 & 1.69 & 1.69 & 1.69 \\
VIF (Wage Interaction) & 1.32 & 1.06 & 1.32 & 1.06 & 1.07 & 1.02 \\
VIF (Price Interaction) & 1.32 & 1.06 & 1.32 & 1.06 & 1.07 & 1.02 \\
\midrule
\multicolumn{7}{l}{\textbf{Cointegration Test Triad}} \\
Resid. ADF ($t_{\text{ADF}}$) [$H_0$: No Coint.] & -4.16$^{*}$ & -4.46$^{**}$ & -4.16$^{*}$ & -4.46$^{**}$ & -4.42$^{**}$ & -4.79$^{**}$ \\
Shin LM ($C_{\text{Shin}}$) [$H_0$: Coint.] & 0.05 & 0.03 & 0.05 & 0.03 & 0.04 & 0.02 \\
VW Type 2 ($S_{\text{Type 2}}$) [$H_0$: Coint.] & 16.11 & 9.55 & 16.11 & 9.55 & 16.40 & 8.61 \\
\midrule
Bandwidth Ratio ($b$) & 0.2088 & 0.1758 & 0.2088 & 0.1758 & 0.2133 & 0.1758 \\
Sample Size ($T$) & 94 & 94 & 94 & 94 & 94 & 94 \\
\bottomrule
\multicolumn{7}{p{15cm}}{\footnotesize \textbf{Note:} Full cointegration diagnostics for the estimates reported in the main text. \textbf{Star reversal warning:} for the Residual ADF test, significance stars indicate rejection of No Cointegration (supporting cointegration). For the Shin LM test, significance stars indicate rejection of Cointegration (indicating spurious regression). The Vogelsang--Wagner Type 2 test shares the Shin null ($H_0$: cointegration) and its stars read the same way. Critical values: Residual ADF (MacKinnon 2010, Case 2, constant only, $m=4$) $-3.81$ / $-4.10$ / $-4.64$ at 10\% / 5\% / 1\%; Shin LM (Shin 1994, Table 1, demeaned $C_\mu$) $0.163$ / $0.221$ / $0.380$ for $k=2$ and $0.121$ / $0.159$ / $0.271$ for $k=3$, where $k$ counts stochastic $I(1)$ regressors; Vogelsang--Wagner Type 2 critical values are fixed-$b$ and jointly indexed by $(b, k)$, so no single constant applies across models \textbf{[PENDING: see the note to the corresponding main-text table]}. Bandwidth ratio $b$ is the model's HAC bandwidth divided by $T$.} \\
\end{tabular}
\end{table}


\begin{table}[H]
\centering
\setlength{\tabcolsep}{2.5pt}
\caption{Specification Set 2: Effective Wage Interaction (GVA Basis) --- Cointegration Diagnostics and Collinearity}
\label{tab:specB_set2_gva_triad}
\scriptsize
\begin{tabular}{p{3.8cm}p{1.1cm}p{1.1cm}p{1.1cm}p{1.1cm}p{1.1cm}p{1.1cm}}
\toprule
\multicolumn{7}{l}{\textbf{Dependent Variable: $y_t^{\text{NFC}}$ (Log NFC Real GVA)}} \\
\midrule
Model & M2.13a & M2.14b & M2.17a & M2.18b & M2.21a & M2.22b \\
\midrule
\multicolumn{7}{l}{\textbf{Collinearity Diagnostics (VIF)}} \\
VIF ($k^{\text{NRC}}$) & 1.69 & 1.69 & 1.69 & 1.69 & 1.69 & 1.69 \\
VIF ($\kappa$) & 1.69 & 1.69 & 1.69 & 1.69 & 1.69 & 1.69 \\
VIF (Wage Interaction) & 1.18 & 1.03 & 1.03 & 1.00 & 1.00 & 1.00 \\
VIF (Price Interaction) & 1.18 & 1.03 & 1.03 & 1.00 & 1.00 & 1.00 \\
\midrule
\multicolumn{7}{l}{\textbf{Cointegration Test Triad}} \\
Resid. ADF ($t_{\text{ADF}}$) [$H_0$: No Coint.] & -4.33$^{*}$ & -4.58$^{**}$ & -4.40$^{*}$ & -4.77$^{**}$ & -4.44$^{**}$ & -4.88$^{**}$ \\
Shin LM ($C_{\text{Shin}}$) [$H_0$: Coint.] & 0.03 & 0.02 & 0.03 & 0.02 & 0.05 & 0.02 \\
VW Type 2 ($S_{\text{Type 2}}$) [$H_0$: Coint.] & 14.97 & 5.93 & 14.58 & 3.71 & 16.05 & 3.21 \\
\midrule
Bandwidth Ratio ($b$) & 0.1854 & 0.1731 & 0.1921 & 0.1735 & 0.2125 & 0.1646 \\
Sample Size ($T$) & 94 & 94 & 94 & 94 & 94 & 94 \\
\bottomrule
\multicolumn{7}{p{15cm}}{\footnotesize \textbf{Note:} Full cointegration diagnostics for the estimates reported in the main text. \textbf{Star reversal warning:} for the Residual ADF test, significance stars indicate rejection of No Cointegration (supporting cointegration). For the Shin LM test, significance stars indicate rejection of Cointegration (indicating spurious regression). The Vogelsang--Wagner Type 2 test shares the Shin null ($H_0$: cointegration) and its stars read the same way. Critical values: Residual ADF (MacKinnon 2010, Case 2, constant only, $m=4$) $-3.81$ / $-4.10$ / $-4.64$ at 10\% / 5\% / 1\%; Shin LM (Shin 1994, Table 1, demeaned $C_\mu$) $0.163$ / $0.221$ / $0.380$ for $k=2$ and $0.121$ / $0.159$ / $0.271$ for $k=3$, where $k$ counts stochastic $I(1)$ regressors; Vogelsang--Wagner Type 2 critical values are fixed-$b$ and jointly indexed by $(b, k)$, so no single constant applies across models \textbf{[PENDING: see the note to the corresponding main-text table]}. Bandwidth ratio $b$ is the model's HAC bandwidth divided by $T$.} \\
\end{tabular}
\end{table}


\begin{table}[H]
\centering
\setlength{\tabcolsep}{2.5pt}
\caption{Specification Set 2: Effective Wage Interaction (NVA Basis) --- Cointegration Diagnostics and Collinearity}
\label{tab:specB_set2_nva_triad}
\scriptsize
\begin{tabular}{p{3.8cm}p{1.1cm}p{1.1cm}p{1.1cm}p{1.1cm}p{1.1cm}p{1.1cm}}
\toprule
\multicolumn{7}{l}{\textbf{Dependent Variable: $y_t^{\text{NFC}}$ (Log NFC Real GVA)}} \\
\midrule
Model & M2.15a & M2.16b & M2.19a & M2.20b & M2.23a & M2.24b \\
\midrule
\multicolumn{7}{l}{\textbf{Collinearity Diagnostics (VIF)}} \\
VIF ($k^{\text{NRC}}$) & 1.69 & 1.69 & 1.69 & 1.69 & 1.69 & 1.69 \\
VIF ($\kappa$) & 1.69 & 1.69 & 1.69 & 1.69 & 1.69 & 1.69 \\
VIF (Wage Interaction) & 1.18 & 1.03 & 1.18 & 1.03 & 1.02 & 1.00 \\
VIF (Price Interaction) & 1.18 & 1.03 & 1.18 & 1.03 & 1.02 & 1.00 \\
\midrule
\multicolumn{7}{l}{\textbf{Cointegration Test Triad}} \\
Resid. ADF ($t_{\text{ADF}}$) [$H_0$: No Coint.] & -4.33$^{*}$ & -4.58$^{**}$ & -4.33$^{*}$ & -4.58$^{**}$ & -4.44$^{**}$ & -4.86$^{**}$ \\
Shin LM ($C_{\text{Shin}}$) [$H_0$: Coint.] & 0.03 & 0.02 & 0.03 & 0.02 & 0.05 & 0.02 \\
VW Type 2 ($S_{\text{Type 2}}$) [$H_0$: Coint.] & 14.97 & 5.93 & 14.97 & 5.93 & 15.72 & 4.59 \\
\midrule
Bandwidth Ratio ($b$) & 0.1854 & 0.1731 & 0.1854 & 0.1731 & 0.2137 & 0.1661 \\
Sample Size ($T$) & 94 & 94 & 94 & 94 & 94 & 94 \\
\bottomrule
\multicolumn{7}{p{15cm}}{\footnotesize \textbf{Note:} Full cointegration diagnostics for the estimates reported in the main text. \textbf{Star reversal warning:} for the Residual ADF test, significance stars indicate rejection of No Cointegration (supporting cointegration). For the Shin LM test, significance stars indicate rejection of Cointegration (indicating spurious regression). The Vogelsang--Wagner Type 2 test shares the Shin null ($H_0$: cointegration) and its stars read the same way. Critical values: Residual ADF (MacKinnon 2010, Case 2, constant only, $m=4$) $-3.81$ / $-4.10$ / $-4.64$ at 10\% / 5\% / 1\%; Shin LM (Shin 1994, Table 1, demeaned $C_\mu$) $0.163$ / $0.221$ / $0.380$ for $k=2$ and $0.121$ / $0.159$ / $0.271$ for $k=3$, where $k$ counts stochastic $I(1)$ regressors; Vogelsang--Wagner Type 2 critical values are fixed-$b$ and jointly indexed by $(b, k)$, so no single constant applies across models \textbf{[PENDING: see the note to the corresponding main-text table]}. Bandwidth ratio $b$ is the model's HAC bandwidth divided by $T$.} \\
\end{tabular}
\end{table}


\subsection{Chile: threshold specification}

\begin{table}[H]
\centering
\setlength{\tabcolsep}{3.5pt}
\caption{Chile Peripheral Extension: Threshold CMPR IM-OLS Cointegration Estimates (1920--2024) --- Cointegration Diagnostics and Collinearity}
\label{tab:chile_peripheral_imols_triad}
\scriptsize
\begin{tabular}{p{4.2cm}p{1.4cm}p{1.4cm}p{1.4cm}p{1.4cm}p{1.4cm}}
\toprule
\multicolumn{6}{l}{\textbf{Dependent Variable: $y_t$ (Log Aggregate Chilean Real GDP)}} \\
\midrule
Parameter / Variable & Spec M3.1a (Baseline) & Spec M3.2a (Mediated) & Spec M3.3a (T-CMPR Slack $E_A$) & Spec M3.3b (T-CMPR Trap $E_B$) & Spec M3.4a (Full T-CMPR) \\
\midrule
\multicolumn{6}{l}{\textbf{Collinearity Diagnostics (VIF)}} \\
VIF ($k^{\text{NRC}}$) & 1.42 & 1.42 & 1.42 & 1.42 & 1.42 \\
VIF ($\kappa$) & 1.42 & 1.42 & 1.42 & 1.42 & 1.42 \\
VIF (Wage / Threshold Interactions) & 1.02 & 1.05 & 1.02 & 1.04 & 1.04 \\
\midrule
\multicolumn{6}{l}{\textbf{Cointegration Test Triad}} \\
Resid. ADF ($t_{\text{ADF}}$) [$H_0$: No Coint.] & -4.12$^{**}$ & -4.35$^{**}$ & -4.40$^{**}$ & -4.18$^{**}$ & -4.62$^{***}$ \\
Shin LM ($C_{\text{Shin}}$) [$H_0$: Coint.] & 0.04 & 0.03 & 0.03 & 0.04 & 0.02 \\
VW Type 2 ($S_{\text{Type 2}}$) [$H_0$: Coint.] & 9.15 & 7.82 & 7.25 & 11.40 & 6.85 \\
\midrule
Sample Size ($T$) & 104 & 104 & 104 & 104 & 104 \\
\bottomrule
\multicolumn{6}{p{11.2cm}}{\footnotesize \textbf{Note:} Full cointegration diagnostics for the estimates reported in the main text. \textbf{Star reversal warning:} for the Residual ADF test, significance stars indicate rejection of No Cointegration (supporting cointegration). For the Shin LM test, significance stars indicate rejection of Cointegration (indicating spurious regression). The Vogelsang--Wagner Type 2 test shares the Shin null ($H_0$: cointegration) and its stars read the same way. Critical values: Residual ADF (MacKinnon 2010, Case 2, constant only, $m=4$) $-3.81$ / $-4.10$ / $-4.64$ at 10\% / 5\% / 1\%; Shin LM (Shin 1994, Table 1, demeaned $C_\mu$) $0.163$ / $0.221$ / $0.380$ for $k=2$ and $0.121$ / $0.159$ / $0.271$ for $k=3$, where $k$ counts stochastic $I(1)$ regressors; Vogelsang--Wagner Type 2 critical values are fixed-$b$ and jointly indexed by $(b, k)$, so no single constant applies across models \textbf{[PENDING: see the note to the corresponding main-text table]}. Bandwidth ratio $b$ is the model's HAC bandwidth divided by $T$.} \\
\end{tabular}
\end{table}


\subsection{The Monte Carlo protocol}

\paragraph{A}
Coefficient-level significance throughout the chapter comes from a 100,000-pass Monte Carlo simulation that reproduces the exact data-generating process of each specification variant, including whether step level-shifts are present. This is not a convenience. Fixed-$b$ critical values for a single-coefficient $t$-test are systematically larger in absolute value than the standard-normal cutoffs of 1.645, 1.960, and 2.576, so applying the normal cutoffs to an estimator whose stated methodology is fixed-$b$ HAR overstates significance. Re-auditing Specification A under the corrected procedure moved 39 coefficient-level stars, every one of them a downgrade.

\paragraph{B}
For the Chilean threshold specification the simulation additionally replicates the Profile Least Squares grid search over the capacity-to-import threshold inside each pass. Without this the inference would be conditional on a threshold treated as known when it was in fact estimated from the same data, and the grid search consumes degrees of freedom that the asymptotic bounds do not account for.

\paragraph{C}
\PLACEHOLDER{Tail-convergence evidence for the simulation --- the MATHPLACEHOLDER0XYZ, MATHPLACEHOLDER1XYZ, and MATHPLACEHOLDER2XYZ convergence diagnostics at MATHPLACEHOLDER3XYZ, and the inflation of the 95th-percentile critical values under structural breaks --- exists at \texttt{chapter2\_vault/06\_paper\_facing/results/03\_monte\_carlo\_tail\_convergence\_report.md} with four supporting figures, and should be folded into this subsection.}
\FloatBarrier
\section{Parameter Stability}

\paragraph{A}
A cointegration test that fails should be diagnosed, not merely reported. This appendix sets out how instability is tested for in this chapter and what constraints govern reading a failure as evidence of regime change rather than as evidence of no long-run relation.

\subsection{What the Shin statistic already tells us}

\paragraph{A}
Shin's $C$ statistic is Hansen's $L_c$ restricted to intercept-only parameter variation, with $L_c^{\text{intercept}} = T \cdot C_{\text{Shin}}$. This identity holds to machine precision across the Specification A models.

\paragraph{B}
The consequence is interpretive and it matters. A Shin rejection is already evidence of parameter instability; it is not evidence of "no cointegration" in some undifferentiated sense. A report that observes a Shin rejection and stops there has not said what it appears to have said. It must either decompose the rejection with the full Hansen $L_c$, per parameter and jointly, or state plainly that the decomposition was not performed.

\paragraph{C}
One numerical caveat travels with this identity wherever both statistics appear: $L_c$ and $C_{\text{Shin}}$ will not agree numerically, because they use different long-run variance estimators — a score-process Bartlett estimator against the conditional $\omega_{u \cdot v}$ from the cointegrating regression. On this data they differ by roughly an order of magnitude. The discrepancy is expected and is not a sign that either is wrong.

\subsection{Bootstrapped nulls, and why the direction of the null matters}

\paragraph{A}
Test statistics introduced into this chapter take their inference from a bootstrapped null distribution rather than from hardcoded critical values, unless a properly indexed published table has been digitized in full.

\paragraph{B}
The construction of that null is where the single most dangerous error in this protocol lives. Tests whose null is cointegration with stable parameters — Hansen $L_c$, MeanF, SupF, and the CUSUM family — require resampling stationary residuals under a moving-block scheme. Tests whose null is no cointegration — Gregory-Hansen — require cumulating resampled differences into a random walk. Transposing the two data-generating processes silently inverts every $p$-value while producing output that looks entirely plausible. Block length follows the fitted long-run-variance bandwidth; i.i.d. resampling is not admissible, since it destroys exactly the serial dependence the block scheme exists to preserve.

\subsection{Which specification the stability tests run on}

\paragraph{A}
Endogenously selected step dummies are excluded from the stability tests. The CUSUM family and Gregory-Hansen run on the core specification — intercept plus stochastic regressors — because a dummy selected from the same data whose stability is being tested contaminates the test. Hansen $L_c$ is reported both on the full fitted design and restricted to core parameters, and any table reporting it states which.

\paragraph{B}
The distinction between level shifts and vector shifts governs what may be claimed. Step dummies absorb level shifts only. They cannot absorb a shift in the cointegrating vector itself. A specification that answers detected breaks with step dummies has therefore not addressed slope instability, and is not described as break-robust anywhere in this chapter on that basis.

\subsection{When a failure may be read as regime change}

\paragraph{A}
Reading a cointegration failure as evidence of structural regime-switching is permitted only when four conditions hold jointly.

\paragraph{B}
Admissibility must not be the explanation: the integration-order gate passes for every regressor as actually constructed, lag structure included. Instability must be affirmatively demonstrated with bootstrapped $p$-values, not inferred from the failure itself. Break dates must be corroborated across the Hansen SupF argmax, the CUSUM supF argmax, Bai-Perron, and Gregory-Hansen — where these converge the reading is strengthened, and where they diverge the divergence is reported as divergence rather than resolved to whichever date suits the argument. And the constructive question must be answered: does a long-run relation appear once the cointegrating vector is allowed to break, by regime-split re-estimation or by Gregory-Hansen? If it does not, the honest verdict is that no stable long-run relation has been identified, and the text says so.

\paragraph{C}
Describing a cointegration failure as evidence of regime change on the strength of demonstrated instability alone, without the other three conditions, is not available.

\subsection{Small samples}

\paragraph{A}
Where a sub-sample falls to 30 observations or fewer, bootstrap $p$-values are reported alongside asymptotics. Below roughly 25 observations the estimate is flagged as underpowered and carries an explicit reliability caveat rather than a bare verdict. Step dummies that are constant within a sub-sample are dropped, and where this happens the text notes that the sub-sample specification therefore differs from the full-sample one.

\paragraph{B}
\PLACEHOLDER{Stability results are not yet reported here. Required: Hansen MATHPLACEHOLDER0XYZ (full design and core-restricted), MeanF, SupF, CUSUM, and Gregory-Hansen for each preferred specification, each with bootstrapped MATHPLACEHOLDER1XYZ-values under the correctly oriented null; the break-date corroboration table across the four detection methods; and the consolidated verdict summary. Source material is in the archived \texttt{reports/US\_Section31\_Stability\_Audit/99\_archive/}.}

\paragraph{C}
\PLACEHOLDER{One known defect to resolve before those results are reported: the residual ADF implementation used in the Chilean structural-break chain is not lag-augmented, and an unaugmented Dickey-Fuller regression on serially correlated residuals under-rejects severely. On the Specification A models the difference is decisive --- 2 of 18 rejections at the 5\% level with zero lagged differences against 16 of 18 under an information-criterion-selected lag order. Any prior reading of cointegration failure as regime-switching that rests on the unaugmented statistic must be revisited before it is repeated here.}

% ---------- Bibliography ----------%
\pagebreak
\bibliographystyle{apalike}
\bibliography{references}

\end{document}
